\section{Technical Lemmas}
\begin{lemma}\label{lem:gaussian-sub-exp}
    If $X \sim \mathcal{N}(\mu, \sigma^2)$ and $Y := X^2 - \mathbb{E}[X^2]$,  
then for all $|\lambda| \leq \frac{1}{4\sigma^2}$, we have:
\[
\log \mathbb{E}(e^{\lambda Y}) \leq 4 \left( \sigma^4 + \mu^2 \sigma^2 \right) \lambda^2
\]
\end{lemma}
\textbf{Proof:} 

Let $X \sim \mathcal{N}(\mu, \sigma^2)$. For $|\lambda| \leq \frac{1}{4\sigma^2}$, we compute:

\[
\mathbb{E}\left[e^{\lambda X^2}\right] 
= \frac{1}{\sqrt{2\pi \sigma^2}} \int_{-\infty}^{\infty} \exp\left( \lambda x^2 - \frac{(x - \mu)^2}{2\sigma^2} \right) dx
\]

Now, complete the square in the exponent:
\[
\lambda x^2 - \frac{(x - \mu)^2}{2\sigma^2} 
= -\frac{1}{2\sigma^2} \left[ (1 - 2\sigma^2 \lambda)x^2 - 2\mu x + \mu^2 \right]
\]

\[
= -\frac{1 - 2\sigma^2 \lambda}{2\sigma^2} 
\left( x-\frac{\mu}{1 - 2\sigma^2 \lambda} \right)^2 
+ \frac{\mu^2 \lambda}{1 - 2\sigma^2 \lambda}
\]

Hence:
\[
\mathbb{E}\left[e^{\lambda X^2}\right] 
= \frac{e^{\mu^2 \lambda / (1 - 2\sigma^2 \lambda)}}{\sqrt{2\pi \sigma^2}} 
\int_{-\infty}^{\infty} 
\exp\left( - \frac{(1 - 2\sigma^2 \lambda)}{2\sigma^2} 
\left( x-\frac{ \mu}{1 - 2\sigma^2 \lambda} \right)^2 \right) dx
\]

This integral is Gaussian:
\[
= \frac{e^{\mu^2 \lambda / (1 - 2\sigma^2 \lambda)}}{\sqrt{2\pi \sigma^2}} 
\cdot \sqrt{ \frac{2\pi \sigma^2}{1 - 2\sigma^2 \lambda} }
\]

\[
= (1 - 2\sigma^2 \lambda)^{-1/2} \exp\left( \frac{\mu^2 \lambda}{1 - 2\sigma^2 \lambda} \right)
\]

As
\[
\lambda \mathbb{E}[X^2] = \lambda(\mu^2 + \sigma^2)
\]

\[
\log \mathbb{E}[e^{\lambda Y}] 
= -\frac{1}{2} \log(1 - 2\sigma^2 \lambda) 
+ \frac{\mu^2 \lambda}{1 - 2\sigma^2 \lambda} 
- \lambda(\mu^2 + \sigma^2)
\]

Let $\nu := 2\sigma^2 \lambda$.  
As $\lambda \leq \frac{1}{4\sigma^2}$, we have $|\nu| \leq \frac{1}{2}$.

\[
\Rightarrow \log \mathbb{E}[e^{\lambda Y}] 
= \left( -\frac{1}{2} \log(1 - \nu) - \frac{\nu}{2} \right) 
+ \mu^2 \lambda \left( \frac{1}{1 - \nu} - 1 \right)
\]

For any $\nu\in[-\tfrac{1}{2}, \tfrac{1}{2}]$, we have:

\begin{enumerate}
    \item 
    \[
    -\frac{1}{2} \log(1 - \nu) - \frac{\nu}{2} \leq \nu^2
    \]

    \item 
    \[
    \left| \frac{1}{1 - \nu} - 1 \right| 
    = \left| \frac{\nu}{1 - \nu} \right| 
    \leq \frac{|\nu|}{1 - |\nu|} 
    \leq 2|\nu|
    \]
\end{enumerate}

Hence we have:
\[
\log \mathbb{E}[e^{\lambda Y}] 
\leq \nu^2 + \mu^2 |\lambda| |2\nu| 
= 4\sigma^4 \lambda^2 + 4\mu^2 \sigma^2 \lambda^2
\]

\begin{lemma}\label{lem:sub-exp-additive}
If $U_1, \ldots, U_K$ are independent mean-zero, sub-exponential random variables with parameters $(V_1, b),\ldots,(V_K,b)$ respectively, then we have the following:
\[
 \Pr\left( \left| \frac{1}{K} \sum_{s=1}^{K} U_s \right| > t \right) 
\leq 2 \exp\left( -cK \min\left\{ \frac{t^2}{\bar{V}}, \frac{t}{b} \right\} \right)
\]
where $\bar{V} = \frac{1}{K} \sum_{k=1}^{K} V_k$ and $c>0$ is some absolute constant.
\end{lemma}
\begin{proof}
Let $U$ be a mean-zero sub-exponential random variables with parameters $(V,b)$. First, we show the following:
    \[
\Pr(|U| > t) \leq 2 \exp\left( -c \min\left\{ \frac{t^2}{V}, \frac{t}{b} \right\} \right)
\]

For any $t \geq 0$ and any $\lambda \in \left[0, \frac{1}{b} \right]$,
\begin{align*}
\Pr[U \geq t] &= \Pr\left[ e^{\lambda U} \geq e^{\lambda t} \right] \\
&\leq e^{-\lambda t} \mathbb{E}[e^{\lambda U}] \\
&\leq \exp\left( -\lambda t + \frac{\lambda^2 V}{2} \right)
\end{align*}

Now we minimize over $\lambda \in \left[0, \frac{1}{b}\right]$.  
Observe that the unconstrained minimizer is: 
\[
\lambda^* = \frac{t}{V}
\]

If $t \leq \frac{V}{b}$, then $\lambda^* \leq \frac{1}{b}$, so we get:
\[
\Pr(U \geq t) \leq \exp\left( -\frac{t^2}{2V} \right)
\]

If $t > \frac{V}{b}$, the minimum occurs at $\lambda = \frac{1}{b}$:
\[
\Pr(U \geq t) \leq \exp\left( -\frac{t}{b} + \frac{V}{2b^2} \right)
\leq \exp\left( -\frac{1}{2} \cdot \frac{t}{b} \right)
\]

If $U_1, \ldots, U_K$ are independent mean-zero, sub-exponential random variables with parameters 
\[
(V_s, b) \quad \text{(same } b \text{ for all)}, 
\]
then for any $|\lambda| \leq \frac{K}{b}$, we have:

\[
\mathbb{E} \left[ \exp\left( \frac{\lambda}{K} \sum_{s=1}^K U_s \right) \right] 
= \prod_{s=1}^K \mathbb{E} \left[ e^{\frac{\lambda}{K} U_s} \right]
\leq \exp\left( \frac{\lambda^2}{2K^2} \sum_{s=1}^K V_s \right)
\]
 \[
\Rightarrow \bar{U} := \frac{1}{K} \sum_{s=1}^K U_s 
\quad \text{is sub-exponential with parameters } 
\left( \frac{1}{K^2} \sum_{s=1}^K V_s, \, \frac{b}{K} \right)
\]

\[
\Rightarrow \Pr\left( \left| \frac{1}{K} \sum_{s=1}^{K} U_s \right| > t \right) 
\leq 2 \exp\left( -cK \min\left\{ \frac{t^2}{\bar{V}}, \frac{t}{b} \right\} \right)
\]

where 
\[
\bar{V} = \frac{1}{K} \sum_{k=1}^{K} V_k
\]
\end{proof}

\begin{lemma}[Hanson–Wright inequality]\label{lem:hanson-wright}
Let $X = (X_1, \ldots, X_n) \in \mathbb{R}^n$ be a random vector with independent components $X_i$ satisfying:
\[
\mathbb{E}[X_i] = 0, \qquad \|X_i\|_{\psi_2} \leq K
\]

Let $A$ be an $n \times n$ matrix. Then for every $t \geq 0$,
\[
\Pr\left( \left| X^\top A X - \mathbb{E}[X^\top A X] \right| > t \right) 
\leq 2 \exp\left( -c \min\left( \frac{t^2}{K^4 \|A\|_F^2}, \, \frac{t}{K^2 \|A\|} \right) \right)
\]

where
\[
\|X\|_{\psi_2} = \sup_{p \geq 1} p^{-1/2} \left( \mathbb{E}|X|^p \right)^{1/p}
\]
\end{lemma}

\begin{lemma}[Chain Rule]\label{kl-chain-rule}
        Let $f(x_1,x_2,\ldots,x_n)$ and $g(x_1,x_2,\ldots,x_n)$ be two joint PDFs for a tuple of random variables $(X_i)_{i\in[n]}$. Let the sample space be $\Omega= \mathbb{R}^{n}$. Then we have the following:
        \begin{equation*}
            KL(f,g)=\int\limits_{\omega\in \Omega}f(\omega)\left(KL(f(X_1),g(X_1))+\sum_{i=2}^n KL(f(X_i|X_{-i}=\omega_{-i}),g(X_i|X_{-i}=\omega_{-i}))\right) d\omega\;
        \end{equation*}
        where $X_{-i}=(X_1,\ldots,X_{i-1})$, $\omega_{-i}=(\omega_1,\ldots,\omega_{i-1})$.
    \end{lemma}
\begin{proof}
    \begin{align*}
            KL(f,g)&=\int\limits_{\omega\in \Omega}f(\omega)\log\left(\frac{f(\omega)}{g(\omega)}\right)\;d\omega\\
            &=\int\limits_{\omega\in \Omega}f(\omega)\log\left(\frac{f(\omega_1)\prod_{i=2}^nf(\omega_i|\omega_{-i})}{g(\omega_1)\prod_{i=2}^ng(\omega_i|\omega_{-i})}\right)\;d\omega\\
            &=\int\limits_{\omega\in \Omega}f(\omega)\left(\log\left(\frac{f(\omega_1)}{g(\omega_1)}\right)+\sum_{i=2}^n\log\left(\frac{f(\omega_i|\omega_{-i})}{g(\omega_i|\omega_{-i})}\right)\right)\;d\omega\\
            &=\int\limits_{\omega\in \Omega}f(\omega)\log\left(\frac{f(\omega_1)}{g(\omega_1)}\right)\;d\omega+\sum_{i=2}^n\int\limits_{\omega\in \Omega}f(\omega)\log\left(\frac{f(\omega_i|\omega_{-i})}{g(\omega_i|\omega_{-i})}\right)\;d\omega\\
            &=\int\limits_{\omega_1\in \mathbb{R}}f(\omega_1)\log\left(\frac{f(\omega_1)}{g(\omega_1)}\right)\;d\omega_1+\sum_{i=2}^n\int\limits_{\omega\in \mathbb{R}^{i}}f(\omega)\log\left(\frac{f(\omega_i|\omega_{-i})}{g(\omega_i|\omega_{-i})}\right)\;d\omega\\
            &=KL(f(X_1),g(X_1))+\sum_{i=2}^n\int\limits_{\omega_{-i}\in \mathbb{R}^{i-1}}f(\omega_{-i})\int\limits_{\omega_{i}\in \mathbb{R}}f(\omega_i|\omega_{-i})\log\left(\frac{f(\omega_i|\omega_{-i})}{g(\omega_i|\omega_{-i})}\right)\;d\omega_id\omega_{-i}\\
            &=KL(f(X_1),g(X_1))+\sum_{i=2}^n\int\limits_{\omega_{-i}\in \mathbb{R}^{i-1}}f(\omega_{-i})KL(f(X_i|X_{-i}=\omega_{-i}),g(X_i|X_{-i}=\omega_{-i}))\;d\omega_{-i}\\
            &=\int_{\omega\in \Omega}f(\omega)KL(f(X_1),g(X_1))\;d\omega+\sum_{i=2}^n\int\limits_{\omega\in \Omega}f(\omega)KL(f(X_i|X_{-i}=\omega_{-i}),g(X_i|X_{-i}=\omega_{-i}))\;d\omega\\
            &=\int\limits_{\omega\in \Omega}f(\omega)\left(KL(f(X_1),g(X_1))+\sum_{i=2}^n KL(f(X_i|X_{-i}=\omega_{-i}),g(X_i|X_{-i}=\omega_{-i}))\right)\;d\omega
    \end{align*}
\end{proof}

\section{Non-Adaptive Fixed Design Bounds}
\subsection{Technical Lemmas for Non-Adaptive Fixed-Design Algorithm}\label{appendix:technical-lemmas}
In this section, we state various technical lemmas and explain how they are used in Section \ref{sec:upper}.
\begin{lemma}[Carath\'eodory's theorem]\label{thm:caratheodory}
Let $S \subset \mathbb{R}^p$ and let $x \in \mathrm{conv}(S)$, where $\mathrm{conv}(S)$ denotes the convex hull of $S$.
Then there exist points $s_0,\dots,s_p \in S$ and coefficients $\alpha_0,\dots,\alpha_p \ge 0$ with
$\sum_{i=0}^p \alpha_i = 1$ such that
\[
x = \sum_{i=0}^p \alpha_i s_i.
\]
In particular, every point in the convex hull of $S$ can be expressed as a convex combination of at most $p+1$ points of $S$.
\end{lemma}
Since, $\mathbb{E}_{\eta\sim\mathcal{N}(0,I_d)}[\max_{x\in\mathcal{X}}\langle x, A(\lambda ;\mathcal{X})^{-1/2}\eta\rangle]$ depends on $\lambda$ through $A(\lambda ;\mathcal{X})=\mathbb{E}_{x\sim\lambda}[xx^\top]$, applying Carath\'eodory's theorem on $S=\{xx^\top:x\in\mathcal{X}\}$ yields a finite supported minimizer $\lambda_1$.

\begin{lemma}[G-optimal design]
    Consider a compact set $\mathcal{X}\subset \mathbb{R}^d$ such that it spans $\mathbb{R}^d$. There exists a distribution $\lambda_*$ over $\mathcal{X}$ such that $\max_{x\in\mathcal{X}}\|x\|^2_{A(\lambda_*;\mathcal{X})^{-1}}=d$ and $|\mathrm{supp}(\lambda_*)|\leq d(d+1)/2$.
\end{lemma}
Applying the above lemma yields a finite supported minimizer $\lambda_2$.


\begin{lemma}[Sudakov--Fernique inequality]\label{thm:sudakov-fernique}
Let $(X_t)_{t\in T}$ and $(Y_t)_{t\in T}$ be two mean-zero Gaussian processes. Assume that for all $t,s\in T$, we have
\[
\mathbb{E}\left[\bigl( X_t - X_s \bigr)^2\right] \le \mathbb{E}\left[\bigl( Y_t - Y_s \bigr)^2\right].
\]
Then
\[
\mathbb{E}\left[\sup_{t\in T} X_t \right]\le \mathbb{E}\left[\sup_{t\in T} Y_t\right].
\]
\end{lemma}

Define $X_x=\langle x, A(\lambda_0;\mathcal X)^{-1/2}\eta\rangle$ and $Y_x=\sqrt2\,\langle x, A(\lambda_1;\mathcal X)^{-1/2}\eta\rangle$ with $\eta\sim\mathcal N(0,I_d)$.  
Then for any $s,t\in\mathcal X$, $X_s-X_t=\langle s-t, A(\lambda_0;\mathcal X)^{-1/2}\eta\rangle$, so by $\E[\eta\eta^\top]=I_d$ we have $\E[(X_s-X_t)^2]=(s-t)^\top A(\lambda_0;\mathcal X)^{-1}(s-t)$ (and similarly $\E[(Y_s-Y_t)^2]=2(s-t)^\top A(\lambda_1;\mathcal X)^{-1}(s-t)$).  
Thus the increment condition of Sudakov--Fernique holds, yielding $\E\sup_{x\in\mathcal X}X_x\le \E\sup_{x\in\mathcal X}Y_x$.



\begin{lemma}[Linear image of a standard Gaussian]\label{lem:linear-gaussian}
Let $A\in\mathbb{R}^{m\times n}$ be a fixed matrix and let $\eta\sim\mathcal{N}(0,I_n)$. Then $A\eta$ is distributionally equivalent to a mean-zero Gaussian vector in $\mathbb{R}^m$ with covariance matrix $AA^\top$. That is $A\eta\stackrel{d}{=}\xi$ where $\xi\sim \mathcal{N}(0,AA^\top)$.
\end{lemma}
We used the above lemma to conclude that $\widehat{\theta}$ is distributionally equivalent to $\theta+(\sum_{t=1}^T x_tx_t^\top)^{-1/2}\eta$ where $\eta\sim\mathcal{N}(0,I_d)$.


\begin{lemma}[Rounding lemma \citep{katz2020empirical,allen2021near}]\label{lem:rounding-single}
Let $\mathcal{Z}\subset \mathbb{R}^d$ be finite with $m:=|\mathcal{Z}|$, and let $N\in\mathbb{N}$. Define
\[
S_N \;:=\;\Bigl\{\kappa\in \mathbb{N}^{m}:\ \sum_{x\in\mathcal{Z}} \kappa_x \le N\Bigr\},
\qquad
C_N \;:=\;\Bigl\{\pi\in [0,N]^{m}:\ \sum_{x\in\mathcal{Z}} \pi_x \le N\Bigr\}.
\]
For any weight vector $v\in \mathbb{R}_+^{m}$, define the (design) matrix
\[
A(v)\;:=\;\sum_{x\in\mathcal{Z}} v_x\, x x^\top \;\in\; \mathbb{S}_d^+.
\]
Let $F:\mathbb{S}_d^+\to \mathbb{R}$ satisfy:
(i) if $A\preceq B$ then $F(A)\ge F(B)$, and
(ii) for all $A\in\mathbb{S}_d^+$ and $t\in(0,1)$, $F(tA)=t^{-1}F(A)$.
Fix $\epsilon\in(0,1/6]$. If $m\ge N\ge 5d/\epsilon^2$, then for every $\pi\in C_N$ there exists an algorithm running in
$\widetilde{O}(m d^2)$ time that outputs $\kappa\in S_N$ such that
\[
F\!\bigl(A(\kappa)\bigr)\;\le\; (1+6\epsilon)\,F\!\bigl(A(\pi)\bigr).
\]
Moreover, for any set $V\subset \mathbb{R}^d$, the functions $F_V,G_V:\mathbb{S}_d^+\to\mathbb{R}$ defined by
\[
F_V(A)\;:=\;\mathbb{E}_{\eta\sim \mathcal{N}(0,I)}\!\left[\max_{v\in V} v^\top A^{-1/2}\eta\right],
\qquad
G_V(A)\;:=\;\max_{v\in V} v^\top A v
\]
satisfy conditions (i) and (ii).
\end{lemma}
For any $T\geq 180d$, we used the above lemma to show the existence of a fixed design $x_1,x_2,\ldots,x_T\in \mathcal{X}$ such that $\tau(A_T)\leq 2\tau(A(\lambda_0;\mathcal{X}))$ where $A_T:=\frac{1}{T}\sum_{i=1}^T x_ix_i^\top$ and $\tau(A):=\mathbb{E}_{\eta\sim\mathcal{N}(0,I_d)}\left[\max_{x\in\mathcal{X}}x^\top A^{-1/2}\eta\right]^2+2\max_{x\in\mathcal{X}}\|x\|_{A^{-1}}^2\log(2/\delta).$ In the lemma, we set $\mathcal{Z}$ as the finite support of $\lambda_0$, $N=T$, and $\pi$ as $\pi_x=T\cdot \lambda_0(x)$.

\begin{lemma}[Borell-TIS inequality]
Let $S \subset \mathbb{R}^d$ be bounded. Let $(V_s)_{s\in S}$ be a Gaussian process such that
\[
\mathbb{E}[V_s] = 0 \quad \text{for all } s \in S.
\]
Define
\[
\sigma^2 = \sup_{s\in S} \mathbb{E}[V_s^2].
\]
Then, for all $u > 0$,
\[
\mathbb{P}\!\left( \left| \sup_{s\in S} V_s - \mathbb{E}\!\sup_{s\in S} V_s \right| \ge u \right)
\le
2 \exp\!\left( -\frac{u^2}{2\sigma^2} \right).
\]
\end{lemma}
We applied Borell-TIS inequality on the gaussian process $V_x=x^\top(\widehat{\theta}-\theta)$. We then used the facts that $V_x$ is distributionally equivalent to $\langle x,(\sum_{t=1}x_tx_t^\top)^{-1/2}\eta\rangle$ where $\eta\sim \mathcal{N}(0,I_d)$ and $ \sup_{x\in \mathcal{X}} \mathbb{E}[V_x^2]=\max_{x\in\mathcal{X}}x^\top(\sum_{t=1}x_tx_t^\top)^{-1}x .$
\subsection{Adaptive Lower Bound}\label{appendix:adaptive-lower-bound}

\begin{theorem}\label{thm:main}
Assume $\mathcal{X}\subset \mathbb{R}^d$ is compact, $\mathrm{span}(\mathcal X)=\mathbb{R}^d$, and $d\geq 2$.
Consider the Gaussian linear observation model
\[
y_t=\langle x_t,\theta\rangle + \eta_t,\qquad \eta_t\sim \mathcal{N}(0,1)\ \text{i.i.d.},
\]
where the (possibly adaptive) actions satisfy $x_t\in\mathcal{X}$.
Any algorithm that outputs $\widehat x\in\mathcal{X}$ satisfying, for all $\theta\in\mathbb{R}^d$,
\[
\langle \widehat x,\theta\rangle \ge \max_{x\in \mathcal X}\langle x,\theta\rangle-\varepsilon
\quad\text{with probability at least } 1-\delta,
\]
must use
\[
n=\Omega\!\left(\frac{d\log(1/\delta)}{\varepsilon^2}\right)
\]
samples for some worst-case $\theta$ and $\delta\in(0,1/16)$.
\end{theorem}

\begin{proof}
Let
\[
B := \mathrm{conv}(\mathcal{X}\cup(-\mathcal{X})).
\]
Let $E$ be the L\"owner ellipsoid of $B$, the minimum-volume ellipsoid containing $B$.
Since $B$ is centrally symmetric and full-dimensional, $E$ is an origin-centered ellipsoid and there exists an invertible linear map $A$ such that
\[
A(E)=\mathbb{B}_2^d
\qquad\text{and hence}\qquad
A(B)\subseteq \mathbb{B}_2^d.
\]
Define
\[
K := A(B),\qquad \mathcal{X}' := A(\mathcal{X}),
\]
so $\mathcal{X}'\subseteq K\subseteq \mathbb{B}_2^d$.

We work in the transformed instance $(\mathcal{X}',\theta')$ where $\theta' := (A^{-1})^\top \theta$.
Indeed, $\langle Ax,\theta'\rangle=\langle x,\theta\rangle$, so any lower bound for $\mathcal{X}'$ implies the same $\mathcal{X}$.
Hence it suffices to prove the statement for $\mathcal{X}'$, and we now drop primes and write $\mathcal{X}\subseteq \mathbb{B}_2^d$ with
\[
K = \mathrm{conv}(\mathcal{X}\cup(-\mathcal{X}))\subseteq \mathbb{B}_2^d,
\]
such that $\mathbb{B}_2^d$ is the L\"owner ellipsoid of $K$.

Since $\mathbb{B}_2^d$ is the minimum-volume ellipsoid containing $K$,
By applying John’s theorem on the polar $K^\circ$ (see Theorem 3.1 from \cite{ball1997elementary}), there exist points $u_1,\dots,u_m \in K\cap \mathbb{S}^{d-1}$ and weights $c_1,\dots,c_m>0$ such that
\begin{equation}\label{eq:john}
\sum_{j=1}^m c_j\, u_j u_j^\top = I_d,
\qquad\text{and hence}\qquad
\sum_{j=1}^m c_j = \mathrm{tr}(I_d)=d.
\end{equation}

We now note that these $u_j$ actually belong to $\mathcal{X}\cup(-\mathcal{X})$.
Fix $j$.
Because $u_j\in \mathbb{S}^{d-1}$ and $K\subseteq \mathbb{B}_2^d$, we have
\[
\max_{x\in K}\langle x,u_j\rangle \le \max_{x\in K}\|x\|_2\|u_j\|_2= 1.
\]
On the other hand, $u_j\in K$ implies $\max_{x\in K}\langle x,u_j\rangle \ge \langle u_j,u_j\rangle =1$.
Hence
\[
\max_{x\in K}\langle x,u_j\rangle = 1.
\]
Since $K=\mathrm{conv}(\mathcal{X}\cup(-\mathcal{X}))$ and the objective is linear, the maximum over $K$ equals the maximum over $\mathcal{X}\cup(-\mathcal{X})$.
Thus there exists $v\in \mathcal{X}\cup(-\mathcal{X})$ such that $\langle v,u_j\rangle=1$.
By Cauchy-Schwarz, $\langle v,u_j\rangle \le \|v\|_2\|u_j\|_2 \le 1$, so equality forces $v=u_j$.
Therefore $u_j\in \mathcal{X}\cup(-\mathcal{X})$.

Define
\[
v_j :=
\begin{cases}
u_j, & u_j\in \mathcal{X},\\
-\,u_j, & u_j\in -\mathcal{X},
\end{cases}
\]
so that $v_j\in \mathcal{X}\cap \mathbb{S}^{d-1}$ for every $j$.
Since $v_j v_j^\top = u_j u_j^\top$, the same weights $c_1,\ldots,c_m$ satisfy
\begin{equation}\label{eq:johnX}
\sum_{j=1}^m c_j\, v_j v_j^\top = I_d,
\qquad\text{and}\qquad
\sum_{j=1}^m c_j = d.
\end{equation}

Let $\mathcal{D}$ be the distribution on $[m]$ such that $\mathbb{P}_{J\sim \mathcal{D}}(J=j)=c_j/d$.
Then \eqref{eq:johnX} implies
\begin{equation}\label{eq:cov}
\mathbb{E}[v_J v_J^\top] = \frac{1}{d}I_d.
\end{equation}

We first prove the following lemma.
\begin{lemma}\label{lem:baseline}
Any algorithm that outputs $\widehat x\in\mathcal{X}\subseteq \mathbb{B}_2^d$ satisfying, for all $\theta\in\mathbb{R}^d$,
\[
\langle \widehat x,\theta\rangle \ge \max_{x\in \mathcal X}\langle x,\theta\rangle-\varepsilon
\quad\text{with probability at least } 1-\delta,
\]
must use
\[
n \ \ge\ \frac{2-\sqrt{2}}{32\,\varepsilon^2}\,\log\!\Big(\frac{1}{4\delta}\Big)
\ =\ \Omega\!\left(\frac{\log(1/\delta)}{\varepsilon^2}\right)
\]
samples for some worst-case $\theta$ and $\delta\in(0,1/16)$.
\end{lemma}
\begin{proof}
Recall that $J\sim\mathcal{D}$ is a random index with $\mathbb{P}_{J\sim \mathcal{D}}(J=j)=c_j/d$.
Hence, we have:
\[
\mathbb E[v_J v_J^\top] = \frac{1}{d}I_d.
\]
Fix $v_1$. Then, we have:
\[
\mathbb E\big[(v_1^\top v_J)^2\big]
= v_1^\top \mathbb E[v_J v_J^\top] v_1
= \frac{1}{d}\|v_1\|_2^2
= \frac{1}{d}.
\]
Hence there exists $k\in[m]$ such that $(v_1^\top v_k)^2\le 1/d$, and thus
$v_1^\top v_k \le 1/\sqrt{d}\le 1/\sqrt{2}$ (since $d\ge 2$).
Therefore
\begin{equation}\label{eq:sep_const}
\|v_1-v_k\|_2^2
=2-2v_1^\top v_k
\ge 2-\sqrt{2}.
\end{equation}
Set $a:=v_1$, $b:=v_k$, and denote $\rho := \|a-b\|_2$. Then \eqref{eq:sep_const} gives $\rho^2\ge 2-\sqrt{2}$.

\noindent
Let $u := (a-b)/\|a-b\|_2 \in \mathbb S^{d-1}$ so that $\langle a,u\rangle-\langle b,u\rangle=\rho$.
Define
\[
\theta_+ := \frac{4\varepsilon}{\rho}\,u,
\qquad
\theta_- := -\frac{4\varepsilon}{\rho}\,u.
\]

Consider an algorithm that uses $n$ samples and outputs $\widehat x$ under each of the hypothesis above while satisfying the conditions in the theorem statement.
Let $p_{\max}:=\max_{x\in\mathcal X}\langle x,u\rangle$ and $p_{\min}:=\min_{x\in\mathcal X}\langle x,u\rangle$.
Note that $p_{\max}\ge \langle a,u\rangle$ and $p_{\min}\le \langle b,u\rangle$.

\noindent
Under $\theta_+$, the optimal value is $\max_{x\in\mathcal X}\langle x,\theta_+\rangle=(4\varepsilon/\rho)\,p_{\max}$.
If $\widehat x$ is $\varepsilon$-optimal for $\theta_+$, then
\[
\frac{4\varepsilon}{\rho}\langle \widehat x,u\rangle
=\langle \widehat x,\theta_+\rangle
\ge \frac{4\varepsilon}{\rho}p_{\max}-\varepsilon
\quad\Rightarrow\quad
\langle \widehat x,u\rangle \ge p_{\max}-\frac{\rho}{4}
\ge \langle a,u\rangle-\frac{\rho}{4}.
\]
Similarly under $\theta_-$, the optimal value is $(4\varepsilon/\rho)(-p_{\min})$ and $\varepsilon$-optimality of $\widehat x$ implies
\[
-\frac{4\varepsilon}{\rho}\langle \widehat x,u\rangle
=\langle \widehat x,\theta_-\rangle
\ge \frac{4\varepsilon}{\rho}(-p_{\min})-\varepsilon
\quad\Rightarrow\quad
\langle \widehat x,u\rangle \le p_{\min}+\frac{\rho}{4}
\le \langle b,u\rangle+\frac{\rho}{4}.
\]
Let
\[
s := \frac{\langle a,u\rangle+\langle b,u\rangle}{2}.
\]
Since $\langle a,u\rangle-\langle b,u\rangle=\rho$, the two implications based on $\varepsilon$-optimality of $\widehat x$ become:
\[
\text{Under }\theta_+:\ \langle \widehat x,u\rangle \ge s+\frac{\rho}{4},
\qquad
\text{Under }\theta_-:\ \langle \widehat x,u\rangle \le s-\frac{\rho}{4},
\]
each with probability at least $1-\delta$.
Therefore the decision rule
\[
\widehat H=
\begin{cases}
\theta_+, & \text{if }\langle \widehat x,u\rangle \ge s,\\
\theta_-, & \text{otherwise},
\end{cases}
\]
distinguishes $\theta_+$ from $\theta_-$ with error at most $\delta$ under each hypothesis.

\noindent
Let $\mathbb P_+$ and $\mathbb P_-$ be the probability laws under $\theta_+$ and $\theta_-$.
By Bretagnolle--Huber inequality, we get
\[
\mathbb P_+(\widehat H=\theta_-)+\mathbb P_-(\widehat H=\theta_+)
\ge \tfrac12 \exp\!\big(-D_{\mathrm{KL}}(\mathbb P_+\|\mathbb P_-)\big).
\]
Since each error is at most $\delta$, the left-hand side is at most $2\delta$, hence
\[
D_{\mathrm{KL}}(\mathbb P_+\|\mathbb P_-)\ge \log\!\Big(\frac{1}{4\delta}\Big).
\]
For Gaussian noise and adaptive actions $x_t\in\mathcal X\subseteq \mathbb B_2^d$,
\begin{align*}
D_{\mathrm{KL}}(\mathbb P_+\|\mathbb P_-)
&= \frac12\sum_{t=1}^n \mathbb E_+\!\big[\langle x_t,\theta_+-\theta_-\rangle^2\big]
= \frac12\sum_{t=1}^n \mathbb E_+\!\Big[\Big\langle x_t,\frac{8\varepsilon}{\rho}u\Big\rangle^2\Big] \\
&\le \frac12\sum_{t=1}^n \Big(\frac{8\varepsilon}{\rho}\Big)^2
= \frac{32n\varepsilon^2}{\rho^2},
\end{align*}
using $|\langle x_t,u\rangle|\le \|x_t\|_2\|u\|_2\le 1$.
Combining the last two inequalities, we have
\[
n \ge \frac{\rho^2}{32\varepsilon^2}\log\!\Big(\frac{1}{4\delta}\Big).
\]
Finally, $\rho^2\ge 2-\sqrt2$ by \eqref{eq:sep_const} yields
\[
n \ge \frac{2-\sqrt2}{32\varepsilon^2}\log\!\Big(\frac{1}{4\delta}\Big).
\]
\end{proof}




Define the alternatives
\[
\theta^{(j)} := 3\varepsilon\, v_j,\qquad j=1,\dots,m,
\]
and consider the hypothesis test
\[
H_0:\ \theta = 0,
\qquad
H_1:\ \theta=\theta^{(J)} \text{ where } J \sim \mathcal{D}.
\]
Let $\mathbb{P}_0$ be the law under $H_0$, $\mathbb{P}_1$ the law under $H_1$, and $\mathbb{P}_{\theta^{(j)}}$ the law under $\theta=\theta^{(j)}$.
For Gaussian noise, convexity of KL and the chain rule yield, for any (possibly adaptive) choice of actions $x_t\in\mathcal{X}$,
\[
D_{\mathrm{KL}}(\mathbb{P}_0\|\mathbb{P}_1)
\le \sum_{j=1}^m \frac{c_j}{d}\,D_{\mathrm{KL}}(\mathbb{P}_0\|\mathbb{P}_{\theta^{(j)}})
= \frac12\sum_{t=1}^N \mathbb{E}\!\Big[\big\langle x_t,\theta^{(J)}\big\rangle^2\Big],
\]
where $N$ is the total number of samples used by the algorithm solving the hypothesis testing problem, and the expectation is under $\mathbb{P}_0$ and the distribution $\mathcal{D}$.
Using \eqref{eq:cov} and $\theta^{(J)}=3\varepsilon v_J$,
\begin{align*}
\mathbb{E}\!\Big[\big\langle x_t,\theta^{(J)}\big\rangle^2\Big]
&= 9\varepsilon^2\, \mathbb{E}\big[\langle x_t,v_J\rangle^2\big]
= 9\varepsilon^2\, x_t^\top \mathbb{E}[v_J v_J^\top] x_t \\
&= 9\varepsilon^2\, x_t^\top\Big(\frac{1}{d}I_d\Big)x_t
= \frac{9\varepsilon^2}{d}\|x_t\|_2^2
\le \frac{9\varepsilon^2}{d}.
\end{align*}
Therefore, we have
\begin{equation}\label{eq:KLmix}
D_{\mathrm{KL}}(\mathbb{P}_0\|\mathbb{P}_1)\le \frac{9N\varepsilon^2}{2d}.
\end{equation}
By Bretagnolle--Huber inequality, we have
\[
\mathbb{P}_0(\widehat{H}=H_1)+\mathbb{P}_1(\widehat{H}=H_0)
\ge \tfrac12\exp\!\big(-D_{\mathrm{KL}}(\mathbb{P}_0\|\mathbb{P}_1)\big)
\ge \tfrac12 \exp\!\left(-\frac{9N\varepsilon^2}{2d}\right).
\]
In particular, if $N \le \frac{2d}{9\varepsilon^2}\log\!\big(\frac{1}{10\delta}\big)$,
then the sum of the two errors exceeds $4\delta$, so at least one error exceeds $2\delta$.
Hence any algorithm with error at most $2\delta$ under both $H_0$ and $H_1$ requires
\begin{equation}\label{eq:testlb}
N=\Omega\!\left(\frac{d\log(1/\delta)}{\varepsilon^2}\right).
\end{equation}

Consider an algorithm $\mathsf{Alg}$ that satisfies the conditions of our theorem statement. We now solve the above hypothesis test using $\mathsf{Alg}$.

Run $\mathsf{Alg}$ for $n$ rounds to obtain $\widehat{x}\in\mathcal{X}$.
Then take
\[
n_{\mathrm{est}} :=  \frac{8\log(2/\delta)}{\varepsilon^2}
\]
additional samples using the constant action $x_t=\widehat{x}$ and let $\widehat{v}$ be the empirical mean of these additional observations. Output $\widehat{H}=H_1$ if $\widehat{v}>\varepsilon$ and $\widehat{H}=H_0$ otherwise.

Gaussian concentration implies
\[
|\widehat{v}-\langle \widehat{x},\theta\rangle|\le \varepsilon/2
\quad\text{with probability at least } 1-\delta.
\]

Under $H_0$, $\langle \widehat{x},\theta\rangle=0$, so $\widehat{v}\le \varepsilon/2$ with probability at least $1-\delta$, hence $\widehat{H}=H_0$ with probability at least $1-\delta$.

Under $H_1$, let us condition on $J=j$.
Since $v_j\in\mathcal{X}\cap\mathbb{S}^{d-1}$,
\[
\max_{x\in\mathcal{X}}\langle x,\theta^{(j)}\rangle \ge \langle v_j,3\varepsilon v_j\rangle = 3\varepsilon.
\]
Therefore $\mathsf{Alg}$ outputs $ \widehat{x}$ such that with probability at least $1-\delta$,
\[
\langle \widehat{x},\theta^{(j)}\rangle \ge 3\varepsilon-\varepsilon = 2\varepsilon.
\]
Therefore, with probability at least $1-2\delta$, we have $\widehat{v}\ge 3\varepsilon/2>\varepsilon$ and therefore $\widehat{H}=H_1$.

Thus we solve the hypothesis testing problem with probability at least $1-2\delta$ under both $H_0$ and $H_1$ using
\[
N = n+n_{\mathrm{est}}
\]
samples.

Finally, Lemma \ref{lem:baseline} implies $n\ge c\,\log(2/\delta)/\varepsilon^2$ for an absolute constant $c>0$.
Hence $n_{\mathrm{est}}\le (8/c)\,n$ and so $N\le (1+8/c)\,n $.
Combining with \eqref{eq:testlb} yields
\[
n=\Omega\!\left(\frac{d\log(1/\delta)}{\varepsilon^2}\right).
\]
\end{proof}



\subsection{Gaussian Width Lower Bound}\label{appendix:gaussian-lower-bound}
Fix $\varepsilon\in (0,1]$, $\delta\in(0,1)$ and an $(\varepsilon,\delta)$-PAC algorithm. Let the algorithm run for $T:=\frac{H_1\log(1/\delta)+H_2}{\varepsilon^2}$ rounds and output $\hat x\in \mathcal{X}$.
Define the simple regret for any $\theta$ as
\[
R_T(\theta)\;:=\;\max_{x\in \mathcal{X}}\langle x,\theta\rangle-\langle \hat x,\theta\rangle \;\ge 0.
\]
Also define the quantity
\[
Z(\theta)\;:=\;\max_{x,y\in \mathcal{X}}\langle x-y,\theta\rangle,
\]
Since $\hat x\in \mathcal{X}$, we have $R_T(\theta)\le Z(\theta)$ for every $\theta$.


As the algorithm is $(\varepsilon,\delta)$-PAC, we have $\mathbb P(R_T(\theta)>\varepsilon)\le \delta$. Hence, taking expectation over $\theta\sim\mathcal N(0,\Sigma)$ where $\Sigma=\tau^2 A^{-1}$, we have
\[
\mathbb{E}_{\theta\sim\mathcal{N}(0,\Sigma)}[R_T(\theta)]\leq \varepsilon+\delta\cdot\mathbb{E}_{\theta\sim\mathcal{N}(0,\Sigma)}[Z(\theta)]=\varepsilon+2\tau\cdot \delta\cdot w(\mathcal{X};A)
\]

Recall that for any $T$ and any algorithm, we have:
\[
\mathbb E_{\theta}[R_T(\theta)]\geq \frac{\tau(1-\tau)}{1+\tau^2} w(\mathcal{X};\mathcal{A})
\]
Hence, using the above two inequalities and the facts that $ w(\mathcal{X};\mathcal{A})\geq w(\mathcal{X})/\sqrt{T}$ and $H_1\leq H_2$, we have the following by setting $\delta=0.1$
\[
\frac{0.12 \cdot w(\mathcal{X})}{\sqrt{T}}\leq  \varepsilon \quad \Rightarrow \quad H_2\geq \frac{0.0144}{\log(10)+1} \cdot w(\mathcal{X})^2
\]
\subsection{Singular Case of the Gaussian Width Lower Bound}\label{appendix:singular-case-lower-bound}
In this section, we prove the following result.
\begin{theorem}
    Assume $\mathcal X\subset\mathbb R^d$ is finite with $\mathrm{span}(\mathcal X)=\mathbb R^d$.
Consider $\delta\in(0,1/2)$ and a fixed design $x_1,\dots,x_T $ such that $\sum_{t=1}^T x_tx_t^\top$ is singular. Then for every (possibly randomized) non-adaptive procedure $\mathcal A$ that outputs $\hat x\in\mathcal{X}$ based on this fixed design, there exists a $\theta\in \mathbb{R}^d$ such that
\[
\mathbb{P}(\max_{x\in\mathcal{X}}\langle x-\hat x,\theta\rangle> \varepsilon)>\delta.
\]
\end{theorem}
\begin{proof}
    Since $A:=\sum_{t=1}^T x_tx_t^\top$ is singular, pick $v\neq 0$ such that $Av=0$. Then we have:
\[
0 \;=\; v^\top Av \;=\; \mathbb \sum_{t=1}^T (v^\top x_t)^2.
\]

Therefore, for all $t\in[T]$, we have:
\[
\langle x_t,v\rangle=0
\]

\medskip
Consider two reward vectors
\[
\theta_{+1}=\alpha v,\qquad \theta_{-1}=-\alpha v,
\]
with $\alpha>0$. Then for every $t$ and $s\in \{-1,+1\}$, we have
\[
y_t=\langle x_t,\theta_s\rangle+\eta_t
=s\cdot\alpha\langle x_t,v\rangle+\eta_t
=\eta_t
\]
Hence the distribution of the algorithm's output $\hat x$ is identical under
$\theta_{+1}$ and $\theta_{-1}$.

Define
\[
a:=\max_{x\in \mathcal X}\langle x,v\rangle,\qquad
b:=\min_{x\in \mathcal X}\langle x,v\rangle.
\]
We claim $a>b$. For the sake of contradiction, assume $a=b$. In this case $\langle x,v\rangle$ is constant over $\mathcal X$. However, $\langle x_t,v\rangle$ and $x_t\in \mathcal{X}$ which implies that $\langle x,v\rangle=0$ for all $x\in\mathcal X$. This implies
$v\perp \mathrm{span}(\mathcal X)$, contradicting $\mathrm{span}(\mathcal X)=\mathbb R^d$.
Thus
\[
\Delta(v)\;:=\;a-b\;>\;0.
\]

Let us define regret as $r(\hat x,\theta)=\max_{x\in\mathcal X}\langle x-\hat x,\theta\rangle$. We now compute regret under each $\theta_s$ for $s\in\{-1,+1\}$. 

\begin{itemize}
\item For $\theta_{+1}=\alpha v$,
\[
r(\hat x,\theta_{+1})
=\max_{x\in\mathcal X}\langle x-\hat x,\alpha v\rangle
=\alpha\Big(a-\langle \hat x,v\rangle\Big),
\]

\item For $\theta_{-1}=-\alpha v$,
\[
r(\hat x,\theta_{-1})
=\max_{x\in\mathcal X}\langle x-\hat x,-\alpha v\rangle
=\alpha\Big(\langle \hat x,v\rangle-b\Big),
\]
\end{itemize}

Adding the two bounds gives
\[
r(\hat x,\theta_{+1})+r(\hat x,\theta_{-1})
\;=\;\alpha(a-b)=\alpha\,\Delta(v).
\]
Therefore, if $\alpha=\frac{4\varepsilon}{\Delta(v)}$, we have
\[
\max\Big\{r(\hat x,\theta_{+1}),r(\hat x,\theta_{-1})\Big\}
\;\ge\;\frac{\alpha\,\Delta(v)}{2}=2\varepsilon.
\]
This implies that $\max_{s\in\{-1,+1\}}\mathbb{P}(\max_{x\in\mathcal{X}}\langle x-\hat x,\theta_s\rangle> \varepsilon)\geq \frac{1}{2}>\delta.$
\end{proof}


\section{Adaptive Version of Our Non-Adaptive Algorithm.}\label{appendix:adaptive-version}
In this section, we make our non-adaptive algorithm from Section \ref{sec:upper} adaptive.
Let $\mathcal{R}=(R_1,R_2,\ldots,R_d)$ be a partition of $\mathcal{X}$ into $d$ regions. Let us define 
\[
w(\lambda,\mathcal{X},\mathcal{R})=\max_{i\in[d]}\mathbb{E}_{\eta\sim\mathcal{N}(0,I_d)}\left[\max_{x\in R_i}\langle x, A(\lambda ;\mathcal{X})^{-1/2}\eta\rangle\right]
\]
Let $i_*\in[d]$ be the index such that $x_*\in R_{i_*}$.

Let $\lambda_1$ be a distribution over $\mathcal{X}$ that minimizes the expression $\mathbb{E}_{\eta\sim\mathcal{N}(0,I_d)}[\max_{x\in\mathcal{X}}\langle x, A(\lambda ;\mathcal{X})^{-1/2}\eta\rangle]$. Similarly let $\lambda_2$ be a distribution over $\mathcal{X}$ that minimizes the expression $\max_{x\in\mathcal{X}}\|x\|^2_{A(\lambda;\mathcal{X})^{-1}}$. Assume that $\lambda_1$ and $\lambda_2$ are finite supported (such minimizers always exist).  Let $\lambda_0$ be a distribution over $\mathcal{X}$ such that we sample from $\lambda_1$ with probability $1/2$ and we sample from $\lambda_2$ with probability $1/2$. 


As $\frac{1}{2}A(\lambda_2;\mathcal{X})\prec A(\lambda_0;\mathcal{X})$, for all $x\in \mathcal{X}$, we have $x^\top A(\lambda_0;\mathcal{X})^{-1} x\leq 2x^\top A(\lambda_2;\mathcal{X})^{-1} x $. As $\lambda_2$ is a G-optimal design, we have $\|x\|^2_{A(\lambda_0;\mathcal{X})^{-1}}\leq 2d$. 

As $\frac{1}{2}A(\lambda_1;\mathcal{X})\prec A(\lambda_0;\mathcal{X})$, due to Sudakov-Fernique inequality, we have
\[
w(\lambda_0,\mathcal{X},\mathcal{R})\leq \sqrt{2}\cdot w(\lambda_1,\mathcal{X},\mathcal{R})\leq \sqrt{2}\cdot w(\mathcal{X}).
\]

Now consider a fixed design $x_1,x_2,\ldots,x_T\in \mathcal{X}$ such that $\tau(A_T,\mathcal{R})\leq 2\tau(A(\lambda_0;\mathcal{X}),\mathcal{R})$ where $A_T:=\frac{1}{T}\sum_{i=1}^T x_ix_i^\top$ and $\tau(A,\mathcal{R}):=\max_{i\in[d]}\mathbb{E}_{\eta\sim\mathcal{N}(0,I_d)}\left[\max_{x\in R_i}x^\top A^{-1/2}\eta\right]^2+2\max_{x\in\mathcal{X}}\|x\|_{A^{-1}}^2\log(4/\delta).$ Such a fixed design exists for any $T\geq 180d$ due to Lemma \ref{lem:rounding-single}. Due to the calculations above, we have $\tau(A_T)\leq 4 w(\lambda_1,\mathcal{X},\mathcal{R})^2+8d\log(4/\delta)$.


Let $y_t=\langle x_t,\theta\rangle+\eta_t$ denote the noisy rewards for our fixed design where $\eta_t\sim\mathcal{N}(0,1)$.
Let $\widehat{\theta}=(\sum_{t=1}^T x_tx_t^\top)^{-1}\sum_{t=1}^T x_ty_t$. Now we observe that $\widehat{\theta}$ is distributionally equivalent to $\theta+(\sum_{t=1}^T x_tx_t^\top)^{-1/2}\eta$ where $\eta\sim\mathcal{N}(0,I_d)$. This enables us to apply Borell-TIS inequality on the gaussian process $V_x=x^\top(\widehat{\theta}-\theta)$ over $R_{i_*}$ and obtain the following with probability $1-\delta/2$:
\begin{equation*}
    \Big|\max_{x\in R_{i_*}}\langle x,\widehat{\theta}-\theta\rangle\Big|\leq \frac{1}{\sqrt{T}}\cdot\mathbb{E}_{\eta\sim\mathcal{N}(0,I_d)}\left[\max_{x\in R_{i_*}}x^\top A_T^{-1/2}\eta\right]+\frac{1}{\sqrt{T}}\cdot\sqrt{2\max_{x\in\mathcal{X}}\|x\|_{A_T^{-1}}^2\log(4/\delta)}\leq \sqrt{\frac{2\tau(A_T)}{T}}
\end{equation*}


Let us compute $x^{(1)},x^{(2)},\ldots,x^{(d)}$ such that $x^{(i)}\in\arg\max_{x\in R_i}\langle x,\widehat{\theta}\rangle$. If we choose $T=1440(\frac{d}{\varepsilon^2}\log(4/\delta)+w(\lambda_1,\mathcal{X},\mathcal{R})^2/\varepsilon^2)$, then $|\max_{x\in R_{i_*}}\langle x,\widehat{\theta}-\theta\rangle|\leq \varepsilon/4$ with probability $1-\delta/2$. This implies that $\langle x^{(i_*)},\theta\rangle\geq \langle x_*,\theta\rangle-\varepsilon/2$ with probability $1-\delta/2$. 

If we run the Median Elimination algorithm from \cite{even2002pac} by treating each $x^{(i)}$ as an arm of a stochastic MAB instance with mean $\langle x^{(i)},\theta\rangle$, we get an index $\hat i$ after additional $\left(\frac{d\log(1/\delta)}{\varepsilon^2}\right)$ samples such that with probability $1-\delta/2$, we have $\max_{i\in[d]}\langle x^{(i)}-x^{(\hat i)},\theta\rangle\leq \varepsilon/2$. We output $x^{(\hat i)}$ as the candidate best arm for $\mathcal{X}$ and due to union bound, we have $\langle x_*-x^{(\hat i)},\theta\rangle\leq \varepsilon$.

In certain cases, if the partition is chosen appropriately, our adaptive algorithm can improve the sample complexity by logarithmic factors. For instance, consider a finite set $\mathcal{X}\subset R^d$. For simplicity of presentation, let us assume that $|\mathcal{X}|$ is multiple of $d$. Let $\mathcal{R}=(R_1,\ldots,R_d)$ be a partition of $\mathcal{X}$ such that for all $i\in[d]$, we have $|R_i|=|\mathcal{X}|/d$. Using Proposition \ref{prop:partition-gaussian-width}, we can show that $w(\lambda_1,\mathcal{X},\mathcal{R})\leq O(\sqrt{d\log(|\mathcal{X}|/d)})$. This implies that our adaptive algorithm has an improved sample complexity of $O\left(\frac{d\log((|\mathcal{X}|/d)/\delta )}{\varepsilon^2}\right)$.

\section{$\ell_2$ Norm Estimation Algorithm}\label{appendix-l2-norm-estimation}
Let $\mathcal{X}$ denote the unit $\ell_2$ ball in $\mathbb{R}^d$, and let $\theta\in\mathbb{R}^d$ be the reward vector. Our goal is to estimate $r:=\|\theta\|_2$ using only samples obtained by querying points in $\mathcal{X}$. In Appendix~\ref{sec:ball-additive-estimate}, we analyze Algorithm~\ref{alg:ball-additive}. In Appendix~\ref{sec:ball-mult-estimate}, we study how to estimate $\|\theta\|_2$ up to a constant multiplicative factor. In Appendix~\ref{sec:ball-large-regime}, we consider the regime where $\|\theta\|_2\geq \sqrt{d}$. Finally, in Appendix~\ref{sec:ball-meta-algo}, we analyze a meta-algorithm that handles all regimes and provides a universal guarantee.

\subsection{Estimating with the Help of a Multiplicative Estimate $r_0$}\label{sec:ball-additive-estimate}
Begin by defining $r:=||\theta||_2$. We now assume that we are given a value $r_0$ such that $\varepsilon<r_0<2\sqrt{d}$ and $\frac{r}{2}< r_0\leq 2r$. In the next section, we show how to satisfy these assumptions. Now we analyse Algorithm \ref{alg:ball-additive}.

In this section we work with sub-exponential random variables. A mean zero $U$ is sub-exponential with parameters $(V, b)$ if
\[
\mathbb{E}[e^{\lambda U}] \leq \exp\left( \frac{\lambda^2 V}{2} \right) \quad \text{for } |\lambda| \leq \frac{1}{b}
\]
If $X \sim \mathcal{N}(\mu, \sigma^2)$ and $Y := X^2 - \mathbb{E}[X^2]$,  
then $Y$ is a sub-exponential with parameters $(V,b)=(8\left( \sigma^4 + \mu^2 \sigma^2 \right),4\sigma^2)$ due to Lemma \ref{lem:gaussian-sub-exp}.



If $U_1, \ldots, U_K$ are independent mean-zero, sub-exponential random variables with parameters $(V_1, b),\ldots,(V_K,b)$ respectively, then we get the following due to Lemma \ref{lem:sub-exp-additive}:
\[
 \Pr\left( \left| \frac{1}{K} \sum_{s=1}^{K} U_s \right| > t \right) 
\leq 2 \exp\left( -cK \min\left\{ \frac{t^2}{\bar{V}}, \frac{t}{b} \right\} \right)
\]
where $\bar{V} = \frac{1}{K} \sum_{k=1}^{K} V_k$ and $c>0$ is some absolute constant.



We now begin with some high-level analysis. Recall $x^{(k)} = \frac{(\varepsilon_1, \ldots, \varepsilon_d)}{\sqrt{d}}$ where $\varepsilon_i\overset{iid}{\sim} \{ \pm 1 \}$ and $ \mu_k = \langle x^{(k)}, \theta \rangle$. Then we have the following:
\[
\mathbb{E}[\mu^2_k] = \theta^\top \mathbb{E}\left[ x^{(k)} (x^{(k)})^\top \right] \theta 
= \theta^\top \left( \frac{1}{d} I_d \right) \theta = \frac{r^2}{d}
\]


Now conditioning on $x^{(k)}$ (hence $\mu_k$), we have
\[
\bar{y}_k \sim \mathcal{N}\left(\mu_k, \frac{1}{s}\right)
\Rightarrow \mathbb{E}[\bar{y}_k^2 \mid \mu_k] = \mu_k^2 + \frac{1}{s}\Rightarrow \mathbb{E}[Z_k \mid \mu_k] = d \mu_k^2
\]




We will now aim to bound
\[
\bar{Z} - r^2 
= \underbrace{\left( \frac{1}{K} \sum_{k=1}^{K} d\mu_k^2 - r^2 \right)}_{\text{Term 1}}
+ \underbrace{\left( \frac{1}{K} \sum_{k=1}^{K} \left( Z_k - d\mu_k^2 \right) \right)}_{\text{Term 2}}
\]

\textbf{Bounding Term 1:}
Note that if $r=0$, then  $\text{Term I}=d \mu_k^2 - r^2=0$. Hence, the non-trivial case is when $r>0$ which we focus on below.
Let $X_k := d \mu_k^2 - r^2$.  
We now show that $X_k$ is sub-exponential with parameters
\[
(V, b) = \left( C r^4, \, C r^2 \right)
\]
for some absolute constant $C$ with the help of Hanson–Wright inequality (see Lemma \ref{lem:hanson-wright}).

Recall that
\[
\mu_k = \left\langle \frac{1}{\sqrt{d}} \varepsilon^{(k)}, \theta \right\rangle,
\qquad \varepsilon_i^{(k)} \overset{iid}{\sim} \{ \pm 1 \}
\]

Now observe that:
\[
d \mu_k^2 = \left( \varepsilon^{(k)} \right)^\top A \varepsilon^{(k)}, \quad \text{where } A := \theta \theta^\top .
\]


As $\mathbb{E}[d \mu_k^2] = r^2$, we apply the Hanson–Wright inequality (we apply Lemma \ref{lem:hanson-wright} with $K = 1$, as Rademacher variables have unit $\psi_2$-norm) to get:
\[
\Pr\left( |X_k| > t \right) 
\leq 2 \exp\left( -c \min\left( \frac{t^2}{\|A\|_F^2}, \frac{t}{\|A\|} \right) \right)
\]

As $A = \theta \theta^\top$, we have
\[
\|A\|_F^2 = \sum_i \theta_i^4 + 2 \sum_{i < j} \theta_i^2 \theta_j^2 
= \left( \sum_i \theta_i^2 \right)^2 = \|\theta\|_2^4 = r^4 
\quad 
\]

Note that $A$ is a rank-1 matrix, and
\[
A \theta = \theta \theta^\top \theta = \|\theta\|_2^2 \theta
\]

Hence, $\theta$ is an eigenvector of $A$ with eigenvalue $\|\theta\|_2^2$.

As the operator norm is equal to the largest eigenvalue, we have
\[
\|A\| = \|\theta\|_2^2 = r^2
\]


Hence, we have:
\[
\Pr\left( |X_k| > t \right) \leq 2 \exp\left( -c \min\left( \frac{t^2}{r^4}, \frac{t}{r^2} \right) \right)
\]



Using the above inequality and applying Lemma \ref{lem:tail-to-sub-exp}, we get that $X_k$ is sub-exponential with parameters $(C r^4, C r^2)$ for some constant $C$.


Hence, we have:
\begin{equation}\label{eq:bounding-term-1}
    \Pr\left( \left| \frac{1}{K} \sum_{k=1}^{K} X_k \right| > t \right) 
\leq 2 \exp\left( -\frac{c}{C}\cdot K \min\left( \frac{t^2}{r^4}, \frac{t}{r^2} \right) \right)
\end{equation}


Choosing $t = \frac{r \varepsilon}{4}$, and $K = c_1 r_0^2 \varepsilon^{-2} \log \frac{4}{\delta}$ for some large constant $c_1$,  
we get:
\[
\left| \frac{1}{K} \sum_{k=1}^{K} d\mu_k^2 - r^2 \right| 
\leq \frac{r \varepsilon}{4} \quad \text{with probability } \geq 1 - \frac{\delta}{4}
\]



\textbf{Bounding Term 2:}  Let $W_k := Z_k - d\mu_k^2 = d\left( \bar{y}_k^2 - \frac{1}{s} - \mu_k^2 \right)$. Conditioning on $\mu_k$, we have:
\[
\sqrt{d} \cdot \bar{y}_k \mid \mu_k \sim \mathcal{N}(\sqrt{d}\cdot\mu_k, \sigma^2), \quad \text{with } \sigma^2 = \frac{d}{s}
\]

Hence $W_k \mid \mu_k$ is sub-exponential with parameters:
\[
V(\mu_k) = 8 \left( \frac{d^2}{s^2} + \mu_k^2 \cdot \frac{d^2}{s} \right), \qquad b = \frac{4d}{s}
\]

Conditioning on $\mu_{1:K} := \mu_1, \ldots, \mu_K$, we get:
\[
\Pr\left( \left| \frac{1}{K} \sum_{k=1}^{K} W_k \right| > t \,\middle|\, \mu_{1:K} \right) 
\leq 2 \exp\left( -cK \min\left( \frac{t^2}{\bar{V}}, \frac{t}{b} \right) \right)
\]

where $
\bar{V} := \frac{1}{K} \sum_{k=1}^{K} V(\mu_k)=8 \left( \frac{d^2}{s^2} + \frac{d^2}{s} \cdot \frac{1}{K} \sum_{k=1}^{K} \mu_k^2 \right).
$ We now aim to upper bound
$
\bar{V}.
$

Recall that $X_k = d \mu_k^2 - r^2$ is sub-exponential with parameters $(Cr^4, Cr^2)$.

Hence, plugging $t=r^2/2$ into Eq. \eqref{eq:bounding-term-1} yields the following, for a sufficiently large constant $c_1$:
\[
\Pr\left( \left| \frac{1}{K} \sum_{k=1}^{K} d \mu_k^2 - r^2 \right| > \frac{r^2}{2} \right) 
\leq \frac{\delta}{8}
\]

%for any $K \geq c_1 \log \frac{4}{\delta}$

\bigskip

Hence, with probability at least $1 - \frac{\delta}{8}$, we have $
\bar{V} \leq \tau := 8 \left( \frac{d^2}{s^2} + \frac{3dr^2}{2s} \right).
$

\begin{align*}
    \Pr\left( \left| \frac{1}{K} \sum_{k=1}^{K} W_k \right| > t \right)&= \Pr\left( \left| \frac{1}{K} \sum_{k=1}^{K} W_k \right| > t ,\, \bar{V} > \tau \right) +\Pr\left( \left| \frac{1}{K} \sum_{k=1}^{K} W_k \right| > t ,\, \bar{V} \leq \tau \right)  \\
    &\leq \Pr(\bar{V} > \tau) + \Pr\left( \left| \frac{1}{K} \sum_{k=1}^{K} W_k \right| > t \,\middle|\, \bar{V} \leq \tau \right)\cdot \Pr(\bar{V} \leq \tau)\\
    &\leq \frac{\delta}{8} + \underbrace{\Pr\left( \left| \frac{1}{K} \sum_{k=1}^{K} W_k \right| > t \,\middle|\, \bar{V} \leq \tau \right)}_{\text{Term III}}
\end{align*}

Now we have:
\begin{equation}\label{appendix-eq:bounding-term-3}
    \text{Term III}=\Pr\left( \left| \frac{1}{K} \sum_{k=1}^{K} W_k \right| > t \,\middle|\, \bar{V} \leq \tau \right) 
\leq 2 \exp\left( -cK \min\left( \frac{t^2}{\tau}, \frac{t}{b} \right) \right)
\end{equation}

As $s = \frac{c_0 d}{r_0^2}$, we have the following for some absolute constant $C'$:
\[
b = \left( \frac{4d}{s} \right) = \left( \frac{4}{c_0} \cdot {r_0^2} \right)
\quad \text{and} \quad 
\tau = 8 \left( \frac{r_0^4}{c_0^2} + \frac{3r^2 r_0^2}{2c_0} \right)
\leq C' r_0^4 
\]

Also, with 
\[
K = c_1 r_0^2 \varepsilon^{-2} \log \frac{4}{\delta}
\quad \text{,} \quad 
t = \frac{r \varepsilon}{4} \quad \text{ and }\quad\frac{r}{2}\leq r_0\leq 2r
\]
we get that the Term III is at most $\frac{\delta}{8}$ after plugging the above parameters into Eq. \eqref{appendix-eq:bounding-term-3}.


Hence, with probability at least $1 - \delta/2$, we have 
$
\left| \bar{Z} - r^2 \right| \leq \frac{r \varepsilon}{2}. 
$

Let us now assume that $\left| \bar{Z} - r^2 \right| \leq \frac{r \varepsilon}{2}$. As $\bar{Z}\geq 0$ we have $\hat{r}=\sqrt{\bar{Z}}$, which implies the following:
    \[
    |\hat{r} - r| = \left| \sqrt{\bar{Z}} - r \right| 
    = \left| \bar{Z} - r^2 \right| \bigg/ \left( \sqrt{\bar{Z}} + r \right)
    \leq \frac{\frac{r \varepsilon}{2}}{r}
    = \frac{\varepsilon}{2} < \varepsilon
    \]






\subsection{Estimating $r$ up to a Constant Factor}\label{sec:ball-mult-estimate}

\begin{algorithm2e}[H]
\caption{Adaptive multi-scale test for $r=\|\theta\|_2$}
\label{alg:multiscale-test}
\LinesNumbered
\DontPrintSemicolon
\SetKwInOut{Input}{Input}
\SetKwInOut{Output}{Output}
\SetKwFunction{Test}{Test}
\SetKwFunction{Stat}{Statistic}

\Input{$\varepsilon\in(0,1]$, $\delta\in(0,1/3)$, dimension $d$, large absolute constants $c_0,c_1>0$}
\Output{An estimate $r_0$}

\BlankLine
\SetKwProg{Fn}{Function}{}{end}
\Fn{\Stat{$t_j,\delta_j$}}{
    $s_j \leftarrow c_0 \dfrac{d}{t_j^2}$\;
    $K_j \leftarrow c_1 \log\!\left(\dfrac{1}{\delta_j}\right)$\;
    \For{$k\leftarrow 1$ \KwTo $K_j$}{
        Draw a Rademacher unit vector
$x^{(k)} = \frac{(\varepsilon_1, \dots, \varepsilon_d)}{\sqrt{d}}
\quad \text{where } \varepsilon_i \overset{iid}{\sim} \{ \pm 1 \}$\;
        \For{$\ell\leftarrow 1$ \KwTo $s_j$}{
            Observe $y_{k,\ell} = \langle x^{(k)},\theta\rangle + \eta_{k,\ell}$ where $\eta_{k,\ell}\sim\mathcal{N}(0,1)$ i.i.d.\;
        }
        $\bar y_k \leftarrow \dfrac{1}{s_j}\sum_{\ell=1}^{s_j} y_{k,\ell}$\;
        $Z_k \leftarrow d\left(\bar y_k^2 - \dfrac{1}{s}\right)$\;
    }
    $U(t_j) \leftarrow \dfrac{1}{K_j}\sum_{k=1}^{K_j} Z_k$\;
    \KwRet{$U(t_j)$}\;
}

\BlankLine
\Fn{\Test{$t_j,\delta_j$}}{
    $U(t_j) \leftarrow$ \Stat{$t_j,\delta_j$}\;
    \eIf{$U(t_j) \ge \dfrac{3}{2}t_j^2$}{
        \KwRet{$H_1$} \tcp*{$H_1$ is the hypothesis that $r \ge 2t_j$}
    }{
        \KwRet{$H_0$} \tcp*{$H_0$ is the hypothesis that $r \le t_j$}
    }
}

\BlankLine
$j \leftarrow 0$\;
\While{$2^j\cdot\varepsilon<2\sqrt{d}$}{
    $t_j \leftarrow 2^j\cdot\varepsilon$\;
    $\delta_j \leftarrow \dfrac{\delta}{2^{j+2}}$\;
    outcome $\leftarrow$ \Test{$t_j,\delta_j$}\;
    \If{outcome $= H_0$}{
        \textbf{break}\;
    }
    $j \leftarrow j+1$\;
}
$r_0 \leftarrow t_j$\;
\KwRet{$r_0$}\;
\end{algorithm2e}

We aim to output $r_0$ such that, with high probability, it satisfies meaningful properties that depend on the value of $r$. We discuss these properties in detail toward the end of this section.


Now we analyse the algorithm. Fix $j\geq 0$. Recall that under $\texttt{Test}(t_j,\delta_j)$, we return the hypothesis  $H_1 \, (r \geq 2t_j)$ if $U(t_j) \geq \frac{3}{2} t_j^2$ otherwise we return the hypothesis $H_0\, (r \leq t_j)$. 

We now compute the probability of error under the hypothesis $H_0\, (r \leq t_j)$. Error under the hypothesis $H_0$ ($r \leq t_j$) means:
\[
U(t_j) \geq \frac{3}{2} t_j^2 
\quad \text{whereas} \quad 
\mathbb{E}[U(t_j)] = r^2 \leq t_j^2
\]

This requires an upward deviation of at least:
\[
U(t_j) - \mathbb{E}[U(t_j)] \geq \Delta_0 := \frac{3}{2} t_j^2 - r^2 \geq \frac{1}{2} t_j^2
\]


Recall:
\[
U(t_j) - r^2 = \left( \frac{1}{K_j} \sum_k d\mu_k^2 - r^2 \right) + \frac{1}{K_j} \sum_k (Z_k - d\mu_k^2)
\]

Let:
\[
a_0 := \frac{1}{8} t_j^2, \quad \nu_0 := \frac{1}{8} t_j^2
\]

Then the error under $H_0$ satisfies:
\[
\left\{ \text{error under } H_0 \right\} 
\subseteq \left\{ \left| \frac{1}{K_j} \sum_k d\mu_k^2 - r^2 \right| >a_0 \right\}
\cup \left\{ \left| \frac{1}{K_j} \sum_k (Z_k - d\mu_k^2) \right| > \nu_0 \right\}
\]

\bigskip

Recall that $X_k := d \mu_k^2 - r^2$ is sub-exponential with parameters
\[
(V, b) = \left( C r^4, \, C r^2 \right)
\]
for some absolute constant $C$. Now observe that 

Hence, we have the following for some absolute constant $c$:
\[
\Pr\left( \left| \frac{1}{K_j} \sum_{k=1}^{K_j} X_k \right| > a_0 \right) 
\leq 2 \exp\left( -cK_j \min\left( \frac{a_0^2}{r^4}, \frac{a_0}{r^2} \right) \right)
\]

Note that:
\[
\min\left\{ \frac{a_0^2}{r^4}, \frac{a_0}{r^2} \right\} =\min\left\{ \frac{t_j^4}{64r^4}, \frac{t_j^2}{8r^2} \right\} \geq \frac{1}{64}
\]

As $K_j = c_1 \log \left( \frac{1}{\delta_j} \right)$, then for a large universal constant $c_1$, we have the following:
\[
\left| \frac{1}{K_j} \sum_{k=1}^{K_j} d\mu_k^2 - r^2 \right| 
\leq a_0 \quad \text{with probability } \geq 1 - \frac{\delta_j}{4}
\]

Let $W_k := Z_k - d\mu_k^2$

Conditioning on $\mu_k$, we have:
\[
\sqrt{d} \cdot \bar{y}_k \mid \mu_k \sim \mathcal{N}(\sqrt{d}\cdot\mu_k, \sigma^2), \quad \text{with } \sigma^2 = \frac{d}{s_j}
\]

Recall $W_k \mid \mu_k$ is sub-exponential with parameters:
\[
V(\mu_k) = 8 \left( \frac{d^2}{s_j^2} + \mu_k^2 \cdot \frac{d^2}{s_j} \right), \qquad b = \frac{4d}{s_j}
\]

and therefore we had the following for some absolute constant $c$:
\[
\Pr\left( \left| \frac{1}{K_j} \sum_{k=1}^{K_j} W_k \right| > \nu_0 \right)\leq \Pr(\bar{V}>\tau) + \underbrace{2 \exp\left( -cK_j \min\left( \frac{\nu_0^2}{\tau}, \frac{\nu_0}{b} \right) \right)}_{\text{Term I}}
\]

where 
\[
\bar{V} := \frac{1}{K_j} \sum_{k=1}^{K_j} V(\mu_k),\quad \nu_0 = \frac{t_j^2}{8}, \quad b = \frac{4t_j^2}{c_0}, \quad \tau := 8 \left( \frac{d^2}{s_j^2} + \frac{3dr^2}{2s_j}\right) \leq C't_j^4
\]
for some absolute constant $C'$.

Recall that $X_k = d \mu_k^2 - r^2$ is sub-exponential with parameters $(Cr^4, Cr^2)$. As $K_j = c_1 \log \left( \frac{1}{\delta_j} \right)$, we have for some absolute constant $c$ and for a large universal constant $c_1$:
\[
\Pr\left( \left| \frac{1}{K_j} \sum_{k=1}^{K_j} d \mu_k^2 - r^2 \right| > \frac{r^2}{2} \right) 
\leq 2 \exp(-cK_j) \leq \frac{\delta_j}{4}
\]
Hence, we have $\Pr(\bar{V}>\tau)\leq\frac{\delta_j}{4}$.



As $ \frac{\nu_0}{b}=\frac{c_0}{32}, 
\frac{\nu_0^2}{\tau} \geq\frac{1}{64C'}, $ and $ K_j = c_1 \log \left( \frac{1}{\delta_j} \right)$, then for a large universal constant $c_1$, we have the following:
\[
\Pr\left( \left| \frac{1}{K_j} \sum_{k=1}^{K_j} W_k \right| > \nu_0 \right)\leq \Pr(\bar{V}>\tau) +\text{Term I}\leq \frac{\delta_j}{4}+\frac{\delta_j}{4}=\frac{\delta_j}{2}.
\]

Hence, we have the following:
\[
\Pr(\text{error under } H_0) \leq \Pr\left( \left| \frac{1}{K_j} \sum_{k=1}^{K_j} X_k \right| > a_0 \right)+\Pr\left( \left| \frac{1}{K_j} \sum_{k=1}^{K_j} W_k \right| > \nu_0 \right)\leq \frac{3}{4} \delta_j
\]


We next compute the probability of error under the hypothesis $H_1$ ($r \geq 2t_j$). Error under the hypothesis $H_1$ ($r \geq 2t_j$) means:

\[
U(t_j) \leq \frac{3}{2} t_j^2 
\quad \text{whereas} \quad 
\mathbb{E}[U(t_j)] = r^2 \geq 4t_j^2
\]

This requires a downward deviation of at least:
\[
\mathbb{E}[U(t_j)] - U(t_j) \geq \Delta_1 := r^2 - \frac{3}{2} t_j^2 \geq \frac{5}{8} r^2
\]

Recall:
\[
U(t_j) - r^2 = \left( \frac{1}{K_j} \sum_k d\mu_k^2 - r^2 \right) + \frac{1}{K_j} \sum_k (Z_k - d\mu_k^2)
\]


Let:
\[
a_1 := \frac{1}{8} r^2, \quad \nu_1 := \frac{1}{8} r^2
\]

Then the error under $H_1$ satisfies:
\[
\left\{ \text{error under } H_1 \right\} 
\subseteq \left\{ \left| \frac{1}{K_j} \sum_k d\mu_k^2 - r^2 \right| >a_1 \right\}
\cup \left\{ \left| \frac{1}{K_j} \sum_k (Z_k - d\mu_k^2) \right| > \nu_1 \right\}
\]


Recall that $X_k := d \mu_k^2 - r^2$ is sub-exponential with parameters
\[
(V, b) = \left( C r^4, \, C r^2 \right)
\]
for some absolute constant $C$. Now observe that 

Hence, we have the following for some absolute constant $c$:
\[
\Pr\left( \left| \frac{1}{K_j} \sum_{k=1}^{K_j} X_k \right| > a_1 \right) 
\leq 2 \exp\left( -cK_j \min\left( \frac{a_1^2}{r^4}, \frac{a_1}{r^2} \right) \right)
\]

Note that:
\[
\min\left\{ \frac{a_1^2}{r^4}, \frac{a_1}{r^2} \right\} =\min\left\{ \frac{r^4}{64r^4}, \frac{r^2}{8r^2} \right\} = \frac{1}{64}
\]

As $K_j = c_1 \log \left( \frac{1}{\delta_j} \right)$, then for a large universal constant $c_1$, we have the following:
\[
\left| \frac{1}{K_j} \sum_{k=1}^{K_j} d\mu_k^2 - r^2 \right| 
\leq a_1 \quad \text{with probability } \geq 1 - \frac{\delta_j}{4}
\]

Let $W_k := Z_k - d\mu_k^2$

Conditioning on $\mu_k$, we have:
\[
\sqrt{d} \cdot \bar{y}_k \mid \mu_k \sim \mathcal{N}(\sqrt{d}\cdot\mu_k, \sigma^2), \quad \text{with } \sigma^2 = \frac{d}{s_j}
\]

Recall $W_k \mid \mu_k$ is sub-exponential with parameters:
\[
V(\mu_k) = 8 \left( \frac{d^2}{s_j^2} + \mu_k^2 \cdot \frac{d^2}{s_j} \right), \qquad b = \frac{4d}{s_j}
\]

and therefore we had the following for some absolute constant $c$:
\[
\Pr\left( \left| \frac{1}{K_j} \sum_{k=1}^{K_j} W_k \right| > \nu_1 \right)\leq \Pr(\bar{V}>\tau) + \underbrace{2 \exp\left( -cK_j \min\left( \frac{\nu_1^2}{\tau}, \frac{\nu_1}{b} \right) \right)}_{\text{Term II}}
\]

where 
\[
\bar{V} := \frac{1}{K_j} \sum_{k=1}^{K_j} V(\mu_k),\quad \nu_1 = \frac{r^2}{8}, \quad b = \frac{4t_j^2}{c_0}\leq\frac{r^2}{c_0}, \quad \tau := 8 \left( \frac{d^2}{s_j^2} + \frac{3dr^2}{2s_j}\right) \leq C'r^4
\]
for some absolute constant $C'$.

Recall that $X_k = d \mu_k^2 - r^2$ is sub-exponential with parameters $(Cr^4, Cr^2)$. As $K_j = c_1 \log \left( \frac{1}{\delta_j} \right)$, we have the following for some absolute constant $c$ and for a large universal constant $c_1$:
\[
\Pr\left( \left| \frac{1}{K_j} \sum_{k=1}^{K_j} d \mu_k^2 - r^2 \right| > \frac{r^2}{2} \right). 
\leq 2 \exp(-cK_j) \leq \frac{\delta_j}{4}
\]
 Hence, we have $\Pr(\bar{V}>\tau)\leq\frac{\delta_j}{4}$.



As $ \frac{\nu_1}{b}\geq\frac{c_0}{8}, 
\frac{\nu_1^2}{\tau} \geq\frac{1}{64C'}, $ and $ K_j = c_1 \log \left( \frac{1}{\delta_j} \right)$, then for a large universal constant $c_1$, we have the following:
\[
\Pr\left( \left| \frac{1}{K_j} \sum_{k=1}^{K_j} W_k \right| > \nu_1 \right)\leq \Pr(\bar{V}>\tau) +\text{Term II}\leq \frac{\delta_j}{4}+\frac{\delta_j}{4}=\frac{\delta_j}{2}.
\]


Hence, we have:
\[
\Pr(\text{error under } H_1) \leq \frac{\delta_j}{2}+\frac{\delta_j}{4}= \frac{3}{4} \delta_j
\]

We now establish the guarantees of our algorithm. We divide it into $4$ cases.

\textbf{Case 1: $r\leq \varepsilon$:} In this case, for $j=0$, we have $r\leq t_j$ which implies that the hypothesis $H_0$ holds. Which implies with probability at least $1-\frac{3\delta}{16}$, we return $r_0=\varepsilon$.

\textbf{Case 2: $\varepsilon<r< \sqrt{d}$:} In this case, for all $j$ such that $2t_j\leq r$, $H_1$ always holds. Consider the index $j_*$ such that $t_{j_*}\leq r<2t_{j_*}$. If we return $r_0=t_{j_*}$, then $r/2<r_0\leq r$. If we instead proceed the index $j_*+1$, then $H_0$ holds and if we terminate and return $r_0=t_{j_*+1}=2t_{j_*}$, then $r<r_0\leq 2r$. Hence, with probability at least $1-\sum_{j=1}^\infty 3\delta_j/4=1-\frac{3\delta}{8}$, we return $r/2<r_0\leq 2r$.

\textbf{Case 3: $\sqrt{d}\leq r<4\sqrt{d}$:} In this case, for all $j$ such that $2t_j\leq r$, $H_1$ always holds. If there exists an index $j_*$ such that $t_{j_*}<2\sqrt{d}$ and $t_{j_*}\leq r<2t_{j_*}$ and if we return $r_0=t_{j_*}$, then $r/2<r_0\leq r$. If we instead proceed to the index $j_*+1$, then $H_0$ holds and if we terminate and return $r_0=t_{j_*+1}=2t_{j_*}$, then $r<r_0\leq 2r$. Hence, in this scenario, with probability at least $1-\sum_{j=1}^\infty 3\delta_j/4=1-\frac{3\delta}{8}$, we return $r/2<r_0\leq 2r$.

On the other hand, if there is no index $j_*$ such that $t_{j_*}<2\sqrt{d}$ and $t_{j_*}\leq r<2t_{j_*}$, then $H_1$ always holds for all $j$ such that $2^j\varepsilon<2\sqrt{d}$. This implies that with probability at least $1-\sum_{j=1}^\infty 3\delta_j/4=1-\frac{3\delta}{8}$, we return $r_0\geq 2\sqrt{d}$.

\textbf{Case 4: $r\geq 4\sqrt{d}$:} In this case, for all $j$ such that $2^j\varepsilon<2\sqrt{d}$, $H_1$ always holds. This implies that with probability at least $1-\sum_{j=1}^\infty 3\delta_j/4=1-\frac{3\delta}{8}$, we return $r_0\geq 2\sqrt{d}$.

Now we establish the sample complexity of our algorithm. The sample complexity is upper bounded as
\[
\sum_{j=0}^{\lceil\log_2(2\sqrt{d}/\varepsilon)\rceil}s_jt_j\leq\mathcal{O}\left(\sum_{j=0}^\infty(d/\varepsilon)^22^{-j}\log(2^{j+2}/\delta)\right)=\mathcal{O}\left(\frac{d\log(1/\delta)}{\varepsilon^2}\right).
\]

% \subsection{Basic estimation to terminate}
% In order for the previous algorithm to terminate in finite steps, we need to first ensure that the $\|\theta_2\|_2$ is upper bounded by a small quantity. If that's not the case, then we directly estimate $\|\theta_2\|_2$ here itself.



\subsection{Estimation in the Large-Norm Regime}\label{sec:ball-large-regime}
Recall that observe
\[
y_t \;=\; \langle x_t,\theta\rangle + \eta_t, \qquad t = 1,\dots,n,
\]
where $\theta \in \mathbb{R}^d$ and the noise variables are i.i.d.\ $\eta_t \sim \mathcal{N}(0,1)$. Let $r := \|\theta\|_2$. In this section we assume the ``large norm'' condition
\begin{equation}
r^2 \;\ge\; d.
\label{eq:large-norm-assumption}
\end{equation}

Our goal is to construct an estimator $\hat r$ such that, for suitable absolute
constant $C > 0$,
\[
n \;\ge\; C\,\frac{d}{\varepsilon^2}\log\frac{1}{\delta}
\quad\text{ and }\quad
\Pr\big(|\hat r - r| > \varepsilon\big) \le \delta
\]
for all $\theta$ satisfying \eqref{eq:large-norm-assumption}. We first describe out algorithm below.


\begin{algorithm2e}[H]
\caption{Algorithm for Large-Norm Regime}
\label{alg:large-regime}
\DontPrintSemicolon
\SetKwInOut{Input}{Input}
\SetKwInOut{Output}{Output}
\LinesNumbered
\Input{dimension $d$, sample size $n$}
\BlankLine
\For{$t \leftarrow 1$ \KwTo $n$}{
    Draw a Rademacher unit vector $x_t = \frac{(\varepsilon_1, \dots, \varepsilon_d)}{\sqrt{d}}\quad \text{where } \varepsilon_i \overset{iid}{\sim} \{ \pm 1 \}$.\;
    Observe $y_t$ from the model $y_t = \langle x_t,\theta\rangle + \eta_t$ with $\eta_t\sim\mathcal{N}(0,1)$ i.i.d.\;
}
Form $X \in \mathbb{R}^{n\times d}$ whose $t$-th row is $x_t^\top$, and $y \leftarrow (y_1,\dots,y_n)^\top$\;

\If{$X^\top X$ is invertible}{
    $\hat\theta \leftarrow (X^\top X)^{-1} X^\top y$\;
    $\Sigma \leftarrow (X^\top X)^{-1}$\;
}
\Else{
$\hat r\gets \sqrt{d}$\;
\KwRet{$\hat r$}\;
}

$\hat r \leftarrow \sqrt{\max\{\|\hat\theta\|_2^2 - \operatorname{tr}(\Sigma),0\}}$\;

\KwRet{$\hat r$}\;
\end{algorithm2e}


We now start analysing our algorithm. Let us consider the case when $X^\top X$ is invertible. We later show that this holds with high probability. Recall that the least-squares estimator is defined as:
\[
\hat\theta := (X^\top X)^{-1} X^\top y.
\]

Define
\[
\Sigma := (X^\top X)^{-1}, \qquad \Delta := \hat\theta - \theta = \Sigma X^\top \eta.
\]

Conditioning on $X$ and using that $\eta \sim \mathcal{N}(0,I_n)$, we have
\begin{equation}
\Delta \mid X \sim \mathcal{N}(0,\Sigma).
\label{eq:Delta-distribution}
\end{equation}


Conditioning on $X$ and using \eqref{eq:Delta-distribution}, we get
\[
\mathbb{E}[\Delta \mid X] = 0,
\qquad
\mathbb{E}[\|\Delta\|_2^2 \mid X] = \operatorname{tr}(\Sigma).
\]

We wish to estimate $r^2 = \|\theta\|_2^2$. We first have the following:
\[
\|\hat\theta\|_2^2
= \|\theta + \Delta\|_2^2
= \|\theta\|_2^2 + 2\theta^\top\Delta + \|\Delta\|_2^2
= r^2 + 2\theta^\top\Delta + \|\Delta\|_2^2.
\]

Define the debiased estimator
\begin{equation}
\widehat{R} := \|\hat\theta\|_2^2 - \operatorname{tr}(\Sigma).
\label{eq:r2-hat-def}
\end{equation}

Then, we have
\[
\mathbb{E}[\widehat{R} \mid X] = r^2,
\qquad
\mathbb{E}[\widehat{R}] = r^2.
\]

Let
\begin{equation}
Z := \widehat{R} - r^2
= 2\theta^\top\Delta + \big(\|\Delta\|_2^2 - \operatorname{tr}(\Sigma)\big).
\label{eq:Z-def}
\end{equation}







We will now derive a tail bound for $Z$ conditional on $X$. Let $g \sim \mathcal{N}(0,I_d)$, and observe that $\Delta$ is distributionally equivalent to $\Sigma^{1/2} g$. Then \eqref{eq:Z-def} becomes distributionally equivalent to
\[
\tilde Z
= 2\theta^\top\Sigma^{1/2} g
 + \big(g^\top\Sigma g - \operatorname{tr}(\Sigma)\big).
\]

Define
\[
L := 2\theta^\top\Sigma^{1/2} g,
\qquad
Q := g^\top\Sigma g - \operatorname{tr}(\Sigma),
\]

so that
\[
\tilde Z = L + Q.
\]


Conditioning on $X$, the random variable $L$ is Gaussian with mean zero and variance
\[
\operatorname{Var}(L \mid X)
= 4\|\Sigma^{1/2}\theta\|_2^2
= 4\theta^\top\Sigma\theta.
\]

Therefore for any $u>0$,
\begin{equation}
\Pr\big(|L|\ge u \,\big|\, X\big)
\le 2\exp\Big(-\frac{u^2}{8\theta^\top\Sigma\theta}\Big).
\label{eq:L-tail}
\end{equation}

We now bound the term $Q$. We use the Hanson--Wright inequality in the Gaussian case:
if $g \sim \mathcal{N}(0,I_d)$ and $A \in \mathbb{R}^{d\times d}$ is symmetric, then
there exists an absolute constant $c_0 > 0$ such that for all $u>0$,
\begin{equation}
\Pr\big(|g^\top A g - \operatorname{tr}(A)| \ge u\big)
\le 2\exp\left(
  -c_0 \min\left\{
      \frac{u^2}{\|A\|_F^2},
      \frac{u}{\|A\|_{\mathrm{op}}}
  \right\}
\right).
\label{eq:HW}
\end{equation}

Applying \eqref{eq:HW} with $A = \Sigma$, and using
$\|\Sigma\|_F^2 = \operatorname{tr}(\Sigma^2)$, we obtain
\begin{equation}
\Pr\big(|Q|\ge u \,\big|\, X\big)
\le 2\exp\left(
  -c_0 \min\left\{
      \frac{u^2}{\operatorname{tr}(\Sigma^2)},
      \frac{u}{\|\Sigma\|_{\mathrm{op}}}
  \right\}
\right).
\label{eq:Q-tail}
\end{equation}


For any $t>0$,
\[
\{|\tilde Z|\ge t\}
= \{|L+Q|\ge t\}
\subseteq \Big\{|L|\ge \frac{t}{2}\Big\} \cup \Big\{|Q|\ge \frac{t}{2}\Big\},
\]

hence, by the union bound and the distributional equivalence between $\tilde Z$ and $Z$, we get
\begin{equation}
\Pr\big(|Z|\ge t \,\big|\, X\big)
\le \Pr\Big(|L|\ge \frac{t}{2} \Big| X\Big)
 +  \Pr\Big(|Q|\ge \frac{t}{2} \Big| X\Big).
\label{eq:Z-union}
\end{equation}

Using \eqref{eq:L-tail} with $u = t/2$, we get
\[
\Pr\Big(|L|\ge \frac{t}{2} \Big| X\Big)
\le 2\exp\Big(-\frac{(t/2)^2}{8\theta^\top\Sigma\theta}\Big)
= 2\exp\Big(-\frac{t^2}{32\theta^\top\Sigma\theta}\Big).
\]

Using \eqref{eq:Q-tail} with $u = t/2$, we get
\[
\Pr\Big(|Q|\ge \frac{t}{2} \Big| X\Big)
\le 2\exp\left(
  -c_0 \min\left\{
      \frac{t^2}{4\operatorname{tr}(\Sigma^2)},
      \frac{t}{2\|\Sigma\|_{\mathrm{op}}}
  \right\}
\right).
\]

Thus \eqref{eq:Z-union} implies
\begin{equation}
\Pr\big(|Z|\ge t \,\big|\, X\big)
\le
2\exp\Big(-\frac{t^2}{32\theta^\top\Sigma\theta}\Big)
+
2\exp\left(
  -c_0 \min\left\{
      \frac{t^2}{4\operatorname{tr}(\Sigma^2)},
      \frac{t}{2\|\Sigma\|_{\mathrm{op}}}
  \right\}
\right).
\label{eq:Z-two-terms}
\end{equation}

Define
\[
A_1 := \frac{t^2}{32\theta^\top\Sigma\theta},\quad
A_2 := c_0 \frac{t^2}{4\operatorname{tr}(\Sigma^2)},\quad
A_3 := c_0 \frac{t}{2\|\Sigma\|_{\mathrm{op}}}.
\]

Then \eqref{eq:Z-two-terms} can be written as
\[
\Pr(|Z|\ge t\mid X)
\le 2e^{-A_1} + 2e^{-\min\{A_2,A_3\}}
\le 4e^{-\min\{A_1,A_2,A_3\}}.
\]

For any $a,b>0$ we have
\[
\min\left\{\frac{t^2}{a},\frac{t^2}{b}\right\}
= \frac{t^2}{\max\{a,b\}}
\ge \frac{t^2}{a+b}.
\]
Applying this with $a = 32\theta^\top\Sigma\theta$ and
$b = \frac{4}{c_0}\operatorname{tr}(\Sigma^2)$, and absorbing constants, we obtain
for some absolute $c>0$:
\begin{equation}
\Pr(|Z|\ge t\mid X)
\le 4\exp\left(
  -c \min\left\{
    \frac{t^2}{4\theta^\top\Sigma\theta + 2\operatorname{tr}(\Sigma^2)},
    \frac{t}{\|\Sigma\|_{\mathrm{op}}}
  \right\}
\right).
\label{eq:Z-HW-type}
\end{equation}


Standard results on sub-Gaussian random matrices from the Section 4.7 of \cite{vershynin2018high} gives the following.
Since $\mathbb{E}[x_t x_t^\top] = I_d/d$ and the rows are i.i.d.\ sub-Gaussian, there exists
an absolute constant $C_0>1$ such that if
\begin{equation}
n \;\ge\; C_0\big(d + \log\tfrac{2}{\delta}\big),
\label{eq:n-rm-condition}
\end{equation}

then with probability at least $1-\delta/2$,
\begin{equation}
\frac{1}{2d}I_d
\;\preceq\;
\frac{X^\top X}{n}
\;\preceq\;
\frac{2}{d}I_d.
\label{eq:XTX-bounds}
\end{equation}

Multiplying \eqref{eq:XTX-bounds} by $n$ and inverting the inequalities yields
\begin{equation}
\frac{d}{2n}I_d \;\preceq\; \Sigma \;\preceq\; \frac{2d}{n}I_d.
\label{eq:Sigma-sandwich}
\end{equation}

From \eqref{eq:Sigma-sandwich} we obtain, on this ``good event'' for $X$,
\begin{align}
\|\Sigma\|_{\mathrm{op}}
&= \lambda_{\max}(\Sigma)
 \le \frac{2d}{n},
\label{eq:Sigma-op-bound}
\\
\theta^\top\Sigma\theta
&\le \|\Sigma\|_{\mathrm{op}} \|\theta\|_2^2
 \le \frac{2d}{n} r^2,
\label{eq:theta-Sigma-theta-bound}
\\
\operatorname{tr}(\Sigma^2)
&= \sum_{i=1}^d \lambda_i^2 \le d \lambda_{\max}^2
 \le d\left(\frac{2d}{n}\right)^2
 = \frac{4d^3}{n^2}.
\label{eq:Sigma2-trace-bound}
\end{align}

We now use \eqref{eq:theta-Sigma-theta-bound} and \eqref{eq:Sigma2-trace-bound} in
\eqref{eq:Z-HW-type}.
On the good event for $X$,
\begin{align*}
4\theta^\top\Sigma\theta + 2\operatorname{tr}(\Sigma^2)
&\le 4 \cdot \frac{2d}{n}r^2 + 2\cdot\frac{4d^3}{n^2} \\
&= \frac{8d}{n}r^2 + \frac{8d^3}{n^2}.
\end{align*}
Using $r^2\ge d$ and $n\ge d$, we have
\[
\frac{d^3}{n^2} \le \frac{d^2}{n} \le \frac{d}{n}r^2,
\]
hence
\[
\frac{8d^3}{n^2} \le 8\,\frac{d}{n}r^2.
\]
Therefore
\begin{equation}
4\theta^\top\Sigma\theta + 2\operatorname{tr}(\Sigma^2)
\;\le\; \frac{8d}{n}r^2 + 8\,\frac{d}{n}r^2
\;=\; 16\,\frac{d}{n}r^2.
\label{eq:denominator-bound}
\end{equation}

Let us choose
\[
t = \varepsilon r.
\]
On the good event,
\[
\min\left\{\frac{t^2}{4\theta^\top\Sigma\theta + 2\operatorname{tr}(\Sigma^2)},\frac{t}{\|\Sigma\|_{\mathrm{op}}}\right\}
\;\ge\;
\frac{\varepsilon^2 r^2}{16\,\frac{d}{n}r^2}
= \frac{n\varepsilon^2}{16d}.
\]
Using \eqref{eq:Z-HW-type} we obtain the following for some absolute constant $c'>0$:
\begin{equation}
\Pr\big(|Z|\ge \varepsilon r \mid X\big)
\le 4\exp\left(-c'\,\frac{n\varepsilon^2}{d}\right)
\quad\text{on the good event for }X.
\label{eq:Z-large-norm-cond}
\end{equation}

If we choose
\[
n \;\ge\; C_1\,\frac{d}{\varepsilon^2}\log\frac{4}{\delta}
\]

for a sufficiently large absolute constant $C_1$, then the right-hand side of
\eqref{eq:Z-large-norm-cond} is at most $\delta/2$.
Combining with \eqref{eq:n-rm-condition} and the fact that the good event for $X$
has probability at least $1-\delta/2$, we conclude that
\begin{equation}
\Pr\big(|\widehat{R} - r^2| \ge \varepsilon r\big) \le \delta.
\label{eq:r2-hat-concentration}
\end{equation}

Recall that our estimator for $r$ is
\begin{equation}
\hat r := \sqrt{\max\{\widehat{R},0\}}.
\label{eq:r-hat-def}
\end{equation}

On the event
\[
|\widehat{R} - r^2| \le \varepsilon r
\]
and assuming $\varepsilon \le r/2$ (which trivially holds for $d\geq 4$), we have
\[
\widehat{R} \ge r^2 - \varepsilon r \ge r^2 - \tfrac{r^2}{2} = \tfrac{r^2}{2} > 0,
\]
so $\hat r = \sqrt{\widehat{R}}$, and therefore
\[
|\hat r - r|
= \big|\sqrt{\widehat{R}} - \sqrt{r^2}\big|
= \frac{|\widehat{R} - r^2|}{\sqrt{\widehat{R}} + r}
\le \frac{|\widehat{R} - r^2|}{r}
\le \varepsilon.
\]

Hence
\[
\big\{|\widehat{R} - r^2| \le \varepsilon r\big\}
\subseteq
\big\{|\hat r - r| \le \varepsilon\big\}.
\]

Combining with \eqref{eq:r2-hat-concentration}, we obtain
\[
\Pr\big(|\hat r - r| > \varepsilon\big)
\le \Pr\big(|\widehat{R} - r^2| > \varepsilon r\big)
\le \delta.
\]

Finally, we needed
\[
n \;\ge\; C_0\big(d + \log\tfrac{2}{\delta}\big)
\quad\text{and}\quad
n \;\ge\; C_1\,\frac{d}{\varepsilon^2}\log\frac{4}{\delta},
\]
so for some absolute constant $C_2>0$,
\[
n \;\ge\; C_2\,\frac{d}{\varepsilon^2}\log\frac{1}{\delta}
\quad\Longrightarrow\quad
\Pr\big(|\hat r - \|\theta\|_2| > \varepsilon\big)\le \delta
\]
for all $\theta$ satisfying $r^2 \ge d$.

\subsection{Meta Algorithm for $\ell_2$-Norm Estimation}\label{sec:ball-meta-algo}
In section, we describe a meta algorithm that uses Algorithms \ref{alg:ball-additive}, \ref{alg:multiscale-test}, and \ref{alg:large-regime}.

\begin{algorithm2e}[H]
\caption{Meta-algorithm to estimate $\|\theta\|_2^2$}
\label{alg:meta-algo}
\DontPrintSemicolon
\SetKwInOut{Input}{Input}
\SetKwInOut{Output}{Output}
\LinesNumbered
\Input{dimension $d$}
\BlankLine
Run the Algorithm \ref{alg:multiscale-test} and receive the estimate $r_0$\;
\If{$r_0=\varepsilon$}{
    $\hat r\gets r_0$\;
    \KwRet{$\hat r$}\;
}
\If{$r_0\geq2\sqrt{d}$}{
    Run the Algorithm \ref{alg:large-regime} and receive the estimate $\hat r$\;
    \KwRet{$\hat r$}\;
}
\If{$\varepsilon<r_0<2\sqrt{d}$}{
    Run the Algorithm \ref{alg:ball-additive} and receive the estimate $\hat r$\;
    \KwRet{$\hat r$}\;
}
\end{algorithm2e}

Now we claim that Algorithm \ref{alg:meta-algo} takes $\mathcal{O}\left(\frac{d\log(1/\delta)}{\varepsilon^2}\right)$ samples and returns an estimate $\hat{r}$ such that with probability at least $1-\delta$ we have $|r-\hat r|\leq \varepsilon$, where $r:=\|\theta\|_2$. 

First, consider the case where $r\leq \varepsilon$. We showed in Appendix \ref{sec:ball-mult-estimate} that with probability at least $1-\delta$, we have $r_0=\varepsilon$. Hence, probability at least $1-\delta$ we have $|r-\hat r|\leq \varepsilon$.

Next, consider the case where $\varepsilon<r< \sqrt{d}$. We showed in Appendix \ref{sec:ball-mult-estimate} that with probability at least $1-\delta/2$, we have $r/2<r_0\leq 2r$ and $r_0<2\sqrt{d}$. Conditioned on this good event, if $r_0=\varepsilon$, we have $\varepsilon<r<2\varepsilon$ and therefore $|\hat r-r|\leq \varepsilon$. On the other hand, conditioned on this good event, if $r_0>\varepsilon$ Algorithm \ref{alg:ball-additive} is run and it returns $\hat r$ such that probability at least $1-\delta/2$ we have $|r-\hat r|\leq \varepsilon$. Due to union bound, we return $\hat r$ such that with probability at least $1-\delta$ we have $|r-\hat r|\leq \varepsilon$.

Next, consider the case where $\sqrt{d}\leq r<4\sqrt{d}$. Due to the analysis in Appendix \ref{sec:ball-mult-estimate}, with probability at least $1-\delta/2$ we either run Algorithm \ref{alg:ball-additive} with an estimate $r_0$ satisfying $r/2<r_0\leq 2r$ or run the Algorithm \ref{alg:large-regime}. In either case, the algorithm being run returns $\hat r$ such that probability at least $1-\delta/2$ we have $|r-\hat r|\leq \varepsilon$. Due to union bound, we return $\hat r$ such that with probability at least $1-\delta$ we have $|r-\hat r|\leq \varepsilon$.

Finally, consider the case where $r\geq 4\sqrt{d}$. Due to the analysis in Appendix \ref{sec:ball-mult-estimate}, with probability at least $1-\delta/2$ we have $r_0\geq 2\sqrt{d}$. Hence, conditioned on this good event, Algorithm \ref{alg:large-regime} is run and it returns $\hat r$ such that probability at least $1-\delta/2$ we have $|r-\hat r|\leq \varepsilon$. Due to union bound, we return $\hat r$ such that with probability at least $1-\delta$ we have $|r-\hat r|\leq \varepsilon$.


The sample complexity guarantee follows as each algorithm run in the meta-algorithms takes $\mathcal{O}\left(\frac{d\log(1/\delta)}{\varepsilon^2}\right)$ samples.



\section{Tail to Sub-Exponential Technical Lemma}
\begin{lemma}\label{lem:tail-to-sub-exp}
    Assume $U$ is mean-zero and for all $t\ge 0$ and some absolute constant $c>0$, we have:
\[
\Pr(|U|>t)\le 2\exp\!\left(-c\min\!\left(\frac{t^2}{V},\frac{t}{b}\right)\right).
\]
Then there exists an absolute constant $C>0$ such that $U$ is sub-exponential with parameters
\[
\bigl(C(V+b^2),\, Cb\bigr),
\]
i.e.
\[
\mathbb E\bigl[e^{\lambda U}\bigr]\le \exp\!\left(\frac{\lambda^2\,C(V+b^2)}{2}\right)
\quad\text{for all }|\lambda|\le \frac{1}{Cb}.
\]
\end{lemma}
\begin{proof}
First, observe that for all $t\ge 0$, we have:
\begin{equation}\label{eq:sub-exp-1}
    \Pr(|U|>t)
\le 2e^{-c t^2/V}+2e^{-c t/b}.
\end{equation}


For any integer $k\ge 2$,
\[
\mathbb E|U|^k
= k\int_0^\infty t^{k-1}\Pr(|U|>t)\,dt,
\]

so by \eqref{eq:sub-exp-1}, we have:
\[
\mathbb E|U|^k
\le 2k\int_0^\infty t^{k-1}e^{-c t^2/V}\,dt
\;+\;2k\int_0^\infty t^{k-1}e^{-c t/b}\,dt.
\]

Using the definition of gamma function, we have the following:
\[
\int_0^\infty t^{k-1}e^{-c t^2/V}\,dt
=\frac12\left(\frac{V}{c}\right)^{k/2}\Gamma\!\left(\frac{k}{2}\right),
\]
\[
\int_0^\infty t^{k-1}e^{-c t/b}\,dt
=\left(\frac{b}{c}\right)^k\Gamma(k).
\]

Hence
\begin{equation}\label{eq:sub-exp-2}
    \mathbb E|U|^k
\le k\left(\frac{V}{c}\right)^{k/2}\Gamma\!\left(\frac{k}{2}\right)
+2k\left(\frac{b}{c}\right)^k\Gamma(k).
\end{equation}

Using the properties of gamma function from \cite{gubner2021gamma}, we have the following for all $x\geq 1$:
\begin{equation}\label{eq:sub-exp-3}
    \Gamma(x)\le 3\,x^{x-\frac12}e^{-x}.
\end{equation}

Apply \eqref{eq:sub-exp-3} to $x=k/2\ge 1$ and $x=k\ge 2$:
\[
\Gamma(k/2)\le 3\left(\frac{k}{2}\right)^{\frac{k}{2}-\frac12}e^{-k/2},
\qquad
\Gamma(k)\le 3\,k^{k-\frac12}e^{-k}.
\]
Plugging into \eqref{eq:sub-exp-2} and taking $k$-th roots gives:
\begin{equation}\label{eq:sub-exp-4}
    (\mathbb E|U|^k)^{1/k}
\le \left[k\left(\frac{V}{c}\right)^{k/2}\Gamma(k/2)\right]^{1/k}
+\left[2k\left(\frac{b}{c}\right)^k\Gamma(k)\right]^{1/k}\le 2\sqrt{\frac{Vk}{c}} \;+\; 2\frac{bk}{c}
\end{equation}

From \eqref{eq:sub-exp-4} we also get, using $(x+y)^k\le 2^{k-1}(x^k+y^k)$,
\begin{equation}\label{eq:sub-exp-5}
    \mathbb E|U|^k
\le 2^{k-1}\Bigl( (2\sqrt{Vk/c})^k + (2bk/c)^k\Bigr)
\end{equation}

Since $\mathbb E U=0$,
\[
\mathbb E e^{\lambda U}
=1+\sum_{k\ge 2}\frac{\lambda^k\mathbb E[U^k]}{k!}
\le 1+\sum_{k\ge 2}\frac{|\lambda|^k\mathbb E|U|^k}{k!}.
\]

Using \eqref{eq:sub-exp-5}, we get
\[
\mathbb E e^{\lambda U}
\le 1+\frac12\sum_{k\ge 2}\frac{(4|\lambda|)^k\Bigl( (V/c)^{k/2}k^{k/2} + (b/c)^k k^k\Bigr)}{k!}
=1+S_1+S_2,
\]
where
\[
S_1:=\frac12\sum_{k\ge 2}\frac{\bigl(4|\lambda|\sqrt{V/c}\bigr)^k\,k^{k/2}}{k!},
\qquad
S_2:=\frac12\sum_{k\ge 2}\frac{\bigl(4|\lambda|b/c\bigr)^k\,k^k}{k!}.
\]

First, we bound the term $1+S_2$. Use $k!\ge (k/e)^k$, so $k^k/k!\le e^k$. Then
\[
S_2
\le \frac12\sum_{k\ge 2}\bigl(4e|\lambda|b/c\bigr)^k.
\]
If
\begin{equation}\label{eq:sub-exp-6}
    |\lambda|\le \frac{c}{8eb},
\end{equation}
then $r:=4e|\lambda|b/c\le 1/2$, so
\[
S_2\le \frac12\sum_{k\ge 2}r^k
=\frac12\cdot \frac{r^2}{1-r}
\le r^2
\le 16e^2\frac{\lambda^2 b^2}{c^2}.
\]
Hence, using $1+x\le e^x$, we get:
\begin{equation}\label{eq:sub-exp-7}
    1+S_2\le \exp\!\left(16e^2\frac{\lambda^2 b^2}{c^2}\right).
\end{equation}
Next, we bound the term $1+S_1$.
Let
\[
a:=4|\lambda|\sqrt{V/c}.
\]
For $k\ge 2$, letting $m=\lfloor k/2\rfloor$, we have:
\[
k! \ge m!\left(\frac{k}{2}\right)^{k-m}.
\]
This implies
\begin{equation}\label{eq:sub-exp-8}
    \frac{k^{k/2}}{k!}
\le \frac{k^{k/2}}{m!(k/2)^{k-m}}
\le \frac{2^{\lceil k/2\rceil}}{m!}
\le \frac{2^{k/2+1}}{m!}.
\end{equation}
Now we consider both case on $k$, one where $k$ is odd and one where $k$ is even:
\begin{itemize}
  \item $k=2m$ with $m\ge 1$: by \eqref{eq:sub-exp-8}, we have
  \[
  \frac{a^{2m}(2m)^{m}}{(2m)!}\le \frac{a^{2m}2^{m}}{m!}.
  \]
  \item $k=2m+1$ with $m\ge 1$: by \eqref{eq:sub-exp-8}, we have
  \[
  \frac{a^{2m+1}(2m+1)^{m+1/2}}{(2m+1)!}
  \le \frac{a^{2m+1}2^{m+1}}{m!}.
  \]
\end{itemize}

Therefore
\begin{equation}\label{eq:sub-exp-9}
    S_1
\le \frac12\sum_{m\ge 1}\frac{(2a^2)^m}{m!}
\;+\; a\sum_{m\ge 1}\frac{(2a^2)^m}{m!}
= \Bigl(a+\frac12\Bigr)\bigl(e^{2a^2}-1\bigr).
\end{equation}


Using the inequalities $(x-\tfrac12)^2+\tfrac14\ge 0$ and $e^x\ge 1+x$, we have the following: 
\begin{equation}\label{eq:sub-exp-10}
    a+\frac12 \le 1+a^2 \le e^{a^2}.
\end{equation}
% \[

% \tag{10}
% \]
% (The first inequality is $(a-\tfrac12)^2+\tfrac14\ge 0$; the second is $e^x\ge 1+x$.)

Combining \eqref{eq:sub-exp-9} and \eqref{eq:sub-exp-10}, we get:
\[
S_1 \le e^{a^2}(e^{2a^2}-1)=e^{3a^2}-e^{a^2}\le e^{3a^2}-1,
\]
so
\begin{equation}\label{eq:sub-exp-11}
    1+S_1 \le e^{3a^2}=\exp\!\left(48\frac{\lambda^2 V}{c}\right).
\end{equation}


Since $S_1,S_2\ge 0$,
\[
1+S_1+S_2 \le (1+S_1)(1+S_2).
\]

Using \eqref{eq:sub-exp-7} and \eqref{eq:sub-exp-11}, whenever \eqref{eq:sub-exp-6} holds, we have
\[
\mathbb E e^{\lambda U}
\le \exp\!\left(48\frac{\lambda^2 V}{c}\right)\,
\exp\!\left(16e^2\frac{\lambda^2 b^2}{c^2}\right)
= \exp\!\left(\frac{\lambda^2}{2}\Bigl(\frac{96}{c}V+\frac{32e^2}{c^2}b^2\Bigr)\right).
\]

Define
\[
b' := \frac{8e}{c}\,b,
\qquad
V' := \frac{96}{c}V+\frac{32e^2}{c^2}b^2.
\]
Then (6) is exactly $|\lambda|\le 1/b'$, and we have shown
\[
\mathbb E\bigl[e^{\lambda U}\bigr] \le \exp\!\left(\frac{\lambda^2 V'}{2}\right)
\quad\text{for all }|\lambda|\le \frac{1}{b'}.
\]
So $U$ is sub-exponential with parameters $(V',b')$.

Finally, if one desires a single constant multiplying both $(V+b^2)$ and $b$, one can take
\[
C \;:=\; \max\!\left\{\frac{96}{c},\,\frac{32e^2}{c^2},\,\frac{8e}{c}\right\},
\]
so that $V' \le C(V+b^2)$ and $b'\le Cb$. This gives exactly the stated form:
\[
\mathbb E\bigl[e^{\lambda U}\bigr] \le \exp\!\left(\frac{\lambda^2\,C(V+b^2)}{2}\right)
\quad\text{for }|\lambda|\le \frac{1}{Cb}.
\]
\end{proof}
\section{Gaussian Width Properties}\label{appendix:gaussian-width-properties}
% One can show that that $\mathbb{E}_{\eta\sim\mathcal{N}(0,I_d)}[\max_{x\in\mathcal{X}}\langle x, A(\lambda_1 ;\mathcal{X})^{-1/2}\eta\rangle]\geq \Omega(d\log d)$. Moreover one can show that $\mathbb{E}_{\eta\sim\mathcal{N}(0,I_d)}[\max_{x\in\mathcal{X}}\langle x, A(\lambda_1 ;\mathcal{X})^{-1/2}\eta\rangle]\leq O(d^2)$. Moreover, if $\mathcal{X}$ is finite, one can show that $\mathbb{E}_{\eta\sim\mathcal{N}(0,I_d)}[\max_{x\in\mathcal{X}}\langle x, A(\lambda_1 ;\mathcal{X})^{-1/2}\eta\rangle]\leq O(d\log|\mathcal{X}|)$.

For a design distribution $\lambda$ over the set $\mathcal{X}$, let
\[
\Sigma(\lambda):=\mathbb E_{x\sim\lambda}[xx^\top],\qquad
f(\lambda):=\mathbb E_{g\sim\mathcal N(0,I_d)}\Big[\sup_{x\in\mathcal X}\langle x,\Sigma(\lambda)^{-1/2}g\rangle\Big].
\]
Equivalently, we have
\[
f(\lambda)=\mathbb E_{g\sim\mathcal N(0,I_d)}\Big[\sup_{x\in\mathcal X}\langle \Sigma(\lambda)^{-1/2}x, g\rangle\Big].
\]
Recall that $w(\mathcal{X}):=\inf_{\lambda \in \triangle_{\mathcal{X}}}f(\lambda)$. Now we prove the following propositions. 

\begin{proposition}
    $w(\mathcal{X})\leq d$.
\end{proposition}
\begin{proof}



Let $\lambda_G$ be a G-optimal design, that is a distribution that minimizes $\sup_{x\in\mathcal X} x^\top \Sigma(\lambda)^{-1}x$.
\[
\lambda_G \in \arg\min_\lambda\ \sup_{x\in\mathcal X} x^\top \Sigma(\lambda)^{-1}x.
\]

For compact $\mathcal X$ with $\mathrm{span}(\mathcal X)=\mathbb R^d$, we have the following:
\[
\sup_{x\in\mathcal X} x^\top \Sigma(\lambda_G)^{-1}x \ \le\ d.
\]

Equivalently,
\[
\sup_{x\in\mathcal X}\|\Sigma(\lambda_G)^{-1/2}x\|_2 \ \le\ \sqrt d.
\]

 For any fixed $g$, we have the following:
\[
\sup_{x\in\mathcal X}\langle x,\Sigma(\lambda_G)^{-1/2}g\rangle
=\sup_{x\in\mathcal X}\langle \Sigma(\lambda_G)^{-1/2}x, g\rangle
\le \Big(\sup_{x\in\mathcal X}\|\Sigma(\lambda_G)^{-1/2}x\|_2\Big)\,\|g\|_2
\le \sqrt d\,\|g\|_2.
\]

Take expectation:
\[
f(\lambda_G)\le \sqrt d\ \mathbb E\|g\|_2\leq \sqrt d\ \sqrt{\mathbb E\|g\|_2^2}=d.
\]

Therefore, we have the following:
\[
 w(\mathcal{X})=\inf_{\lambda\in\triangle_{\mathcal{X}}} f(\lambda) \le f(\lambda_G) \le d.\
\]
\end{proof}

\begin{proposition}
    For a finite set $\mathcal{X}$ with $|\mathcal X|=m$, we have $w(\mathcal{X})\leq O(\sqrt{d\log m})$.
\end{proposition}
\begin{proof}
Define $u_x:=\Sigma(\lambda_G)^{-1/2}x$. Then $\|u_x\|_2^2\le d$ for all $x\in\mathcal{X}$.

If $g\sim \mathcal{N}(0,I_d)$, then for each $x\in\mathcal{X}$, $Z_x:=\langle u_x,g\rangle$ is Gaussian with variance $\|u_x\|_2^2\le d$, so for $t\ge 0$,
\[
\Pr(Z_x\ge t)\le \exp\!\left(-\frac{t^2}{2d}\right).
\]

Now by applying the union bound, we get the following:
\[
\Pr\Big(\max_{x\in\mathcal X} Z_x \ge t\Big)\le m\exp\!\left(-\frac{t^2}{2d}\right).
\]

Using the fact $\mathbb E[\max_{x\in\mathcal{X}} Z_x]=\int_0^\infty \Pr(\max_{x\in\mathcal{X}} Z_x\ge t)\,dt$ and choosing $t_0=\sqrt{2d\log m}$, we get the following:
\[
\mathbb E[\max_{x\in\mathcal{X}} Z_x]\le t_0+\int_{t_0}^\infty m e^{-t^2/(2d)}\,dt\leq t_0+m\cdot \frac{d}{t_0} e^{-t_0^2/(2d)}
= t_0+ \frac{d}{t_0}\leq O(\sqrt{d\log m}).
\]

Therefore, we have the following:
\[
 w(\mathcal{X})=\inf_{\lambda\in\triangle_{\mathcal{X}}} f(\lambda) \le f(\lambda_G) \le O(\sqrt{d\log |\mathcal{X}|}).\
\]
\end{proof}

Next, recall the definition of $\lambda_1$, $\mathcal{R}$, $w(\lambda_1,\mathcal{R},\mathcal{X})$ for the finite set $\mathcal{X}$ in Appendix \ref{appendix:adaptive-version}.
\begin{proposition}\label{prop:partition-gaussian-width}
    For a finite set $\mathcal{X}$ with $|\mathcal X|=m\geq d$, we have $w(\lambda_1,\mathcal{R},\mathcal{X})\leq O(\sqrt{d\log(m/d)})$.
\end{proposition}
\begin{proof}
For simplicity of presentation, assume that $m$ is a multiple of $d$. Fix an index $i\in [d]$. Define $u_x:=\Sigma(\lambda_G)^{-1/2}x$. Then $\|u_x\|_2^2\le d$ for all $x\in\mathcal{X}$.

If $g\sim \mathcal{N}(0,I_d)$, then for each $x\in\mathcal{X}$, $Z_x:=\langle u_x,g\rangle$ is Gaussian with variance $\|u_x\|_2^2\le d$, so for $t\ge 0$,
\[
\Pr(Z_x\ge t)\le \exp\!\left(-\frac{t^2}{2d}\right).
\]

Now by applying the union bound, we get the following:
\[
\Pr\Big(\max_{x\in R_i} Z_x \ge t\Big)\le (m/d)\exp\!\left(-\frac{t^2}{2d}\right).
\]

Using the fact $\mathbb E[\max_{x\in R_i} Z_x]=\int_0^\infty \Pr(\max_{x\in R_i} Z_x\ge t)\,dt$ and choosing $t_0=\sqrt{2d\log (m/d)}$, we get the following:
\[
\mathbb E[\max_{x\in R_i} Z_x]\le t_0+\int_{t_0}^\infty (m/d) e^{-t^2/(2d)}\,dt\leq t_0+(m/d)\cdot \frac{d}{t_0} e^{-t_0^2/(2d)}
= t_0+ \frac{d}{t_0}\leq O(\sqrt{d\log (m/d)}).
\]

Therefore, we have the following:
\[
w(\lambda_1,\mathcal{R},\mathcal{X}) \le \max_{i\in[d]}\mathbb E[\max_{x\in R_i} Z_x] \le O(\sqrt{d\log(m/d)}).\
\]
\end{proof}





\begin{proposition}
    $w(\mathcal{X})\geq \Omega(\sqrt{d\log d})$.
\end{proposition}
\begin{proof}
Fix any $\lambda$ with $\Sigma:=\Sigma(\lambda)\succ 0$.
Define the set
\[
T := \{\Sigma^{-1/2}x:\ x\in\mathcal X\}\subset\mathbb R^d,
\]
and define the Gaussian process $Z_t:=\langle g,t\rangle$ for $g\sim\mathcal N(0,I_d)$ over $T$.
Then we have:
\[
f(\lambda)=\mathbb E\Big[\sup_{t\in T} Z_t\Big].
\]
Let $U:=\Sigma^{-1/2}X$ where $X\sim\lambda$. Since $\mathbb E[XX^\top]=\Sigma$,
\[
\mathbb E[UU^\top]=\Sigma^{-1/2}\,\mathbb E[XX^\top]\,\Sigma^{-1/2}=I_d.
\]


Let $m:=\lfloor d/2\rfloor$ and $r:=\sqrt{d/2}$.
We now claim that there exists $t_1,\dots,t_m\in T$ such that $\|t_i-t_j\|_2\ge r$ for all $i\ne j$.

We construct them greedily. First pick $t_1$ arbitrarily. Now assume that $t_1,\dots,t_k$ have been chosen for some $1\le k<m$.
Let $V_k:=\mathrm{span}\{t_1,\dots,t_k\}$, let $P_k$ be the orthogonal projection onto $V_k$,
and let $Q_k:=I-P_k$ be the orthogonal projection onto $V_k^\perp$.
Then $Q_k$ is symmetric and idempotent ($Q_k^2=Q_k$). Moreover, for an orthogonal projection, the trace equals the dimension of the subspace it projects onto. Note that $Q_k$ projects onto $V_k^\perp$, which has dimension $d-\dim(V_k)$. Therefore
\[
\mathrm{tr}(Q_k)=\dim(V_k^\perp)=d-\dim(V_k)\ge d-k,
\]
since $\dim(V_k)\le k$.

Since $Q_k$ is an orthogonal projection,
\[
\|Q_k U\|_2^2 = U^\top Q_k^\top Q_k U=U^\top Q_k^2 U= U^\top Q_k U.
\]
Using $\mathrm{tr}(ab)=ab$ for scalars and cyclicity of trace, we have the following:
\[
\mathbb E\|Q_k U\|_2^2
= \mathbb E[U^\top Q_k U]
= \mathbb E[\mathrm{tr}(Q_k U U^\top)]
= \mathrm{tr}\big(Q_k \mathbb E[UU^\top]\big)
= \mathrm{tr}(Q_k I_d)
= \mathrm{tr}(Q_k)
\ge d-k.
\]
Therefore there exists a point $t_{k+1}\in T$ with
\[
\|Q_k t_{k+1}\|_2^2 \ge d-k.
\]
For any $i\le k$, we have $t_i\in V_k$, which implies $Q_k t_i=0$. As orthogonal projections cannot increase the norm, we have the following:
\[
\|t_{k+1}-t_i\|_2 \ge \|Q_k(t_{k+1}-t_i)\|_2
= \|Q_k t_{k+1}\|_2
\ge \sqrt{d-k}
\ge \sqrt{d/2}=r.
\]
This completes the proof of our claim.


Let $S:=\{t_1,\dots,t_m\}$. Recall that $\|t_i-t_j\|_2\ge r$ for all $i\neq j$.
Now define $Z_i:=\langle g,t_i\rangle$ where $g\sim \mathcal{N}(0,1)$.
Then $(Z_i)_{i=1}^m$ is a centered Gaussian vector and for $i\ne j$, we have:
\[
\mathbb E[(Z_i-Z_j)^2] = \|t_i-t_j\|_2^2 \ge r^2.
\]
Let $\xi_1,\dots,\xi_m$ be i.i.d.\ $\mathcal N(0,1)$ and set $Y_i:=(r/\sqrt{2})\xi_i$.
Then for $i\ne j$, we have:
\[
\mathbb E[(Y_i-Y_j)^2] = \frac{r^2}{2}\mathbb E[(\xi_i-\xi_j)^2]=\frac{r^2}{2}\cdot 2=r^2.
\]
Due to Sudakov–Fernique inequality, we have the following:
\[
\mathbb E\Big[\sup_{t\in T} Z_t\Big]\geq\mathbb E\Big[\max_{1\le i\le m} Z_i\Big]
\ge
\mathbb E\Big[\max_{1\le i\le m} Y_i\Big]
=
\frac{r}{\sqrt{2}}\,\mathbb E\Big[\max_{1\le i\le m}\xi_i\Big]\geq \Omega(\sqrt{d\log d}).
\]

This implies that $w(\mathcal{X})\geq \Omega(\sqrt{d\log d})$.
\end{proof}

\begin{proposition}
    When $\mathcal{X}$ is $\{-1,+1\}^d$, $\{0,+1\}^d$ or a unit $\ell_2$-ball, we have $w(\mathcal{X})\geq\Omega(d)$. For $m\leq d/21$, we have $w(\mathcal{X})\geq \Omega(\sqrt{md})$ when $\mathcal{X}$ is an $m$-set.
\end{proposition}
\begin{proof}
    Recall that in Section \ref{sec:upper}, we described an algorithm with an upper bound of $\mathcal{O}\left(\frac{d\log(1/\delta)}{\varepsilon^2}+\frac{w(\mathcal{X})^2}{\varepsilon^2}\right).$
    When $\mathcal{X}$ is $\{-1,+1\}^d$, $\{0,+1\}^d$ or a unit $\ell_2$-ball, our adaptive lower bound results in Table \ref{table:1} imply that $w(\mathcal{X})\geq \Omega(d)$. Similarly when $\mathcal{X}$ is an $m$-set, our adaptive lower bound results in Table \ref{table:1} imply that $w(\mathcal{X})\geq \Omega(\sqrt{md})$ for $m\leq d/21$.
\end{proof}
%\textcolor{red}{Add a small proposition on the claims on the gaussian width of the structured sets.}

\section{Gaussian Width Calculations for Multi-Task MAB}\label{appendix:gaussian-width-multi-task-MAB}
Recall that \(d=\sum_{j=1}^m d_j\). Consider the blocks \(B_j:=\{d_{1:j-1}+1,\dots,d_{1:j}\}\).  Recall that the action set is
\begin{equation}\label{eq:action-set}
\mathcal{X}:=\Big\{x\in\{0,1\}^d:\ \forall j\in[m],\ \sum_{i\in B_j}x_i=1\Big\}.\nonumber
\end{equation}

Let $r := d-m+1$. Since $\mathcal X$ is $r$-dimensional, it is contained in an $r$-dimensional linear subspace
$ \operatorname{span}(\mathcal X)\subseteq \mathbb R^d$.
Let $U\in\mathbb R^{d\times r}$ be a matrix whose columns form an orthonormal basis of $\operatorname{span}(\mathcal X)$, that is $U^\top U = I_r$ and $\operatorname{range}(U)=\operatorname{span}(\mathcal X)$.
Define the coordinate representation of $\mathcal X$ in $\mathbb R^r$ by
\[
\mathcal X_r := \{\,U^\top x : x\in \mathcal X\,\}\subseteq \mathbb R^r.
\]
Then the maps
\[
x \mapsto U^\top x \quad\text{and}\quad z \mapsto Uz
\]
are inverse to each other on $\mathcal X$ and $\mathcal X_r$, respectively. In particular, for every $x\in \mathcal X$,
\[
x = U(U^\top x),
\]
so $U^\top x$ is the unique coordinate vector of $x$ w.r.t the orthonormal basis $\{U_1,\ldots,U_r\}$ where $U_j$ is the $j$-th column of $U$.


We now construct one such matrix $U$. Define the vector $\mu\in\mathbb R^d$ by setting, for each $i\in\{1,\ldots,d\}$,
\[
\mu_i := \frac{1}{d_j}\quad \text{where $j$ is the unique index such that } i\in B_j.
\]
Then, we have:
\[
S := \|\mu\|_2^2 = \sum_{i=1}^d \mu_i^2 = \sum_{j=1}^m d_j\left(\frac{1}{d_j}\right)^2
= \sum_{j=1}^m \frac{1}{d_j},
\]
and we define the unit vector:
\[
u_0 := \frac{\mu}{\|\mu\|_2} = \frac{\mu}{\sqrt S}.
\]
We now define a matrix \(H_n\in\mathbb R^{n\times(n-1)}\). For \(n\ge 2\), define \((H_n)_{r,k}\) for \(\ell\in[n]\), \(k\in[n-1]\) by
\begin{equation}\label{eq:helmert-def}
(H_n)_{\ell,k}=
\begin{cases}
\frac{1}{\sqrt{k(k+1)}} & 1\le \ell\le k,\\[4pt]
-\frac{k}{\sqrt{k(k+1)}} & \ell=k+1,\\[4pt]
0 & \ell\ge k+2.\nonumber
\end{cases}
\end{equation}
Then, we have:
\begin{equation}\label{eq:helmert-props}
H_n^\top H_n=I_{n-1},\qquad H_n^\top \mathbf 1_n=\mathbf{0}_{n-1},
\qquad
H_n H_n^\top = I_n-\frac1n\mathbf 1_n\mathbf 1_n^\top.
\end{equation}
Let \(s_j:=d_{1:j-1}\). We define the entries of the matrix \(Q_j\in\mathbb R^{d\times(d_j-1)}\) as:
\begin{equation}\label{eq:Qj-embed}
(Q_j)_{i,k}=
\begin{cases}
(H_{d_j})_{\,i-s_j,\ k} & i\in B_j,\\
0 & i\notin B_j,
\end{cases}
\qquad k\in[d_j-1].
\end{equation}
Now we define:
\begin{equation}\label{eq:U-def}
U:=\big[u_0\ \ Q_1\ \cdots\ Q_m\big]\in\mathbb R^{d\times r},\qquad
r:=1+\sum_{j=1}^m(d_j-1)=d-m+1.\nonumber
\end{equation}
As \(Q_j,Q_\ell\) have disjoint supports  for \(j\neq \ell\) and due to Eq. \eqref{eq:helmert-props}, we get:
\begin{equation}\label{eq:Q-orth}
Q_j^\top Q_j = H_{d_j}^\top H_{d_j}=I_{d_j-1},\qquad
Q_j^\top Q_\ell=0\ (j\neq \ell).\nonumber
\end{equation}
For any matrix $A$, let $A_{i,:}$ denote the $i$-th row of the matrix $A$. Since \(u_0\) is constant on block \(B_j\), we have:
\begin{align*}\label{eq:u0-Qj-orth}
u_0^\top Q_j
&=\sum_{i\in B_j}(u_0)_i (Q_j)_{i,:}
=\frac{1/d_j}{\sqrt S}\sum_{r=1}^{d_j}(H_{d_j})_{r,:}
=\frac{1/d_j}{\sqrt S}\,\mathbf 1_{d_j}^\top H_{d_j}
=\mathbf{0}_{n-1}.
\end{align*}
Hence
\begin{equation}\label{eq:U-orth}
U^\top U = I_r.\nonumber
\end{equation}
Next, we prove that \(\operatorname{range}(U)=\operatorname{span}(\mathcal X)\). Towards that, define
\begin{equation}\label{eq:Uj-def}
R_j:=\Big\{v\in\mathbb R^d:\operatorname{supp}(v)\subseteq B_j,\ \mathbf 1_{B_j}^\top v=0\Big\}.\nonumber
\end{equation}
where $\mathbf 1_{B_j}$ is the indicator vector of the block $B_j$. From \eqref{eq:helmert-props}–\eqref{eq:Qj-embed}, we get:
\begin{equation}\label{eq:range-Qj}
\operatorname{range}(Q_j)=R_j.\nonumber
\end{equation}
Also \(\mu\in\operatorname{span}(\mathcal{X})\) and $\forall j,\ \forall k\in[d_j],
e_{s_j+k}-e_{s_j+1} \in \operatorname{span}(\mathcal{X})$, so \(R_j\subseteq \operatorname{span}(\mathcal X)\) (since \(\{e_{s_j+k}-e_{s_j+1}\}_{k=2}^{d_j}\) spans \(R_j\)). Thus
\begin{equation}\label{eq:span-inclusion-1}
\operatorname{span}(\mu)\oplus\Big(\bigoplus_{j=1}^m R_j\Big) \subseteq \operatorname{span}(\mathcal X).\nonumber
\end{equation}
where $\oplus$ means direct sum of vector spaces.

Conversely, for any \(x\in \mathcal X\), we have:
\begin{equation}\label{eq:span-inclusion-2}
\forall j\in[m]:\ \mathbf 1_{B_j}^\top(x-\mu)=1-1=0
\quad\Rightarrow\quad x-\mu\in \bigoplus_{j=1}^m R_j,
\quad\Rightarrow\quad x\in \operatorname{span}(\mu)\oplus\Big(\bigoplus_{j=1}^m R_j\Big).\nonumber
\end{equation}
Hence
\begin{equation}\label{eq:spanX-rangeU}
\operatorname{span}(\mathcal X)=\operatorname{span}(\mu)\oplus\Big(\bigoplus_{j=1}^m R_j\Big)
=\operatorname{range}(U),\text{ and indeed } \dim(\operatorname{span}(\mathcal X))=r=d-m+1.\nonumber
\end{equation}
Now the gaussian width term $w(\mathcal{X})$ is defined using the coordinate representation of $\mathcal X$ in $\mathbb R^r$ as follows:
\[
w(\mathcal{X}):=\inf_{\lambda\in\triangle_{\mathcal{X}_r}}\mathbb E_{g_r\sim \mathcal{N}(0,I_r)}\Big[\sup_{x\in \mathcal{X}}\ \langle U^\top x,\ A(\lambda;\mathcal{X}_r)^{-1/2}g_r\rangle\Big]
\]
Fix a probability distribution $\lambda$ on $\mathcal X_r$ such that, if $Z\sim\lambda$ and we map to $\mathcal X$ via
\[
X := UZ \in \mathcal X \subseteq \mathbb R^d,
\]
then the following holds: for each block $B_k$, exactly one coordinate index $j_k \in B_k$ is selected uniformly at random, and the selections $\{j_k\}_{k=1}^m$ are independent across blocks. In other words, viewed in $\mathbb R^d$, $X$ is distributed as a vector obtained by choosing one coordinate uniformly within each block, independently across blocks.  Now define:
\begin{equation}\label{eq:sigma-def}
\Sigma_r:=\mathbb E_{x\sim\lambda}[xx^\top],\nonumber
\end{equation}
where $\Sigma_r\in \mathbb{R}^{r\times r}$. Now the Gaussian width of $\mathcal{X}$ is upper bounded as follows:
\begin{equation}\label{eq:gaussian-width}
w(\mathcal{X})\leq w(\mathcal{X}_r;\Sigma_r):=\mathbb E_{g_r\sim \mathcal{N}(0,I_r)}\Big[\sup_{x\in \mathcal{X}}\ \langle U^\top x,\ \Sigma_r^{-1/2}g_r\rangle\Big]\nonumber
\end{equation}
Observe that \(\Sigma_r := U^\top \Sigma U\) where $\Sigma:=\mathbb{E}_{x\sim\lambda}[Ux(Ux)^\top]$. Due to the definition of $\lambda$, for \(i\in B_j\), \(k\in B_\ell\), we have:
\begin{equation}\label{eq:Sigma-entries}
\Sigma_{ik}=
\begin{cases}
\frac1{d_j}\mathbf 1\{i=k\} & j=\ell,\\[3pt]
\frac1{d_j d_\ell} & j\neq \ell.
\end{cases}\nonumber
\end{equation}
Let \(v\in R_j\) (\(\mathbf 1_{B_j}^\top v=0\), support in \(B_j\)). For \(i\in B_j\),
\begin{align*}\label{eq:Sigma-on-Uj}
(\Sigma v)_i
&=\sum_{k\in B_j}\frac1{d_j}\mathbf 1\{i=k\}v_k
+\sum_{\ell\neq j}\sum_{k\in B_\ell}\frac1{d_j d_\ell}v_k
=\frac1{d_j}v_i+0,
\end{align*}
so \(\Sigma v=\frac1{d_j}v\). Therefore
\begin{equation}\label{eq:Q-blocks}
Q_j^\top \Sigma Q_j=\frac1{d_j}Q_j^\top Q_j=\frac1{d_j}I_{d_j-1},
\qquad
Q_j^\top \Sigma Q_\ell=0\ (j\neq \ell).\nonumber
\end{equation}
For \(i\in B_j\),
\begin{align*}\label{eq:Sigma-mu}
(\Sigma\mu)_i
&=\sum_{k\in B_j}\frac1{d_j}\mathbf 1\{i=k\}\frac1{d_j}
+\sum_{\ell\neq j}\sum_{k\in B_\ell}\frac1{d_j d_\ell}\frac1{d_\ell}\\
&=\frac1{d_j^2}+\sum_{\ell\neq j}\frac1{d_j d_\ell}
=\frac1{d_j}\sum_{\ell=1}^m\frac1{d_\ell}
=S\,\mu_i.
\end{align*}
Thus \(\Sigma\mu=S\mu\) and, since \(\|u_0\|=1\),
\begin{equation}\label{eq:u0-eigs}
u_0^\top\Sigma u_0 = S,\qquad u_0^\top\Sigma Q_j = 0.\nonumber
\end{equation}
Hence, in the explicit basis \(U\),
\begin{equation}\label{eq:Sigma_r}
\Sigma_r
:=U^\top\Sigma U
=
\operatorname{diag}\Big(S,\ \frac1{d_1}I_{d_1-1},\ \dots,\ \frac1{d_m}I_{d_m-1}\Big),
\quad
\Sigma_r^{-1/2}
=
\operatorname{diag}\Big(S^{-1/2},\ \sqrt{d_1}\,I_{d_1-1},\ \dots,\ \sqrt{d_m}\,I_{d_m-1}\Big).
\end{equation}

\bigskip

\subsection*{D) Gaussian Width in \(r=d-m+1\) Dimensions and Evaluation}

Let \(g_r\sim\mathcal N(0,I_r)\), write \(g_r=(g_0,g^{(1)},\dots,g^{(m)})\) with
\(g^{(j)}\sim\mathcal N(0,I_{d_j-1})\) independent. Recall that
\begin{equation}\label{eq:width-intrinsic}
w(\mathcal X)\leq\mathbb E_{g_r\sim\mathcal{N}(0,I_r)}\Big[\sup_{x\in \mathcal X}\ \langle U^\top x,\ \Sigma_r^{-1/2}g_r\rangle\Big].\nonumber
\end{equation}
Also, we have
\begin{equation}\label{eq:u0-x}
\langle u_0,x\rangle=\frac{\langle \mu,x\rangle}{\sqrt S}
=\frac{\sum_{j=1}^m \frac1{d_j}}{\sqrt S}
=\sqrt S,\qquad \forall x\in \mathcal X,\nonumber
\end{equation}
so, from \eqref{eq:Sigma_r}, we have
\begin{align*}
\langle U^\top x,\Sigma_r^{-1/2}g_r\rangle
&=
S^{-1/2}g_0\langle u_0,x\rangle
+\sum_{j=1}^m \sqrt{d_j}\,\langle g^{(j)},Q_j^\top x\rangle\\
&=
g_0+\sum_{j=1}^m \sqrt{d_j}\,\langle g^{(j)},Q_j^\top x\rangle.
\end{align*}
Therefore
\begin{align}\label{eq:width-reduce}
w(\mathcal X)\leq
\mathbb E[g_0]
+\sum_{j=1}^m \sqrt{d_j}\ \mathbb E\Big[\max_{k\in[d_j]}\ \langle g^{(j)},H_{d_j}^\top e_k\rangle\Big]
&=
\sum_{j=1}^m \sqrt{d_j}\ \mathbb E\Big[\max_{k\in[d_j]}\ \langle g^{(j)},H_{d_j}^\top e_k\rangle\Big].
\end{align}
Let \(h^{(j)}:=H_{d_j}g^{(j)}\in\mathbb R^{d_j}\). Then we have
\begin{equation}\label{eq:h-coords}
\langle g^{(j)},H_{d_j}^\top e_k\rangle = e_k^\top H_{d_j}g^{(j)} = (h^{(j)})_k,\nonumber
\end{equation}
and by \eqref{eq:helmert-props},
\begin{equation}\label{eq:h-law}
h^{(j)}\sim\mathcal N\Big(0,\ H_{d_j}H_{d_j}^\top\Big)
=\mathcal N\Big(0,\ I_{d_j}-\frac1{d_j}\mathbf 1\mathbf 1^\top\Big).\nonumber
\end{equation}
If \(Z^{(j)}\sim\mathcal N(0,I_{d_j})\) and \(\bar Z^{(j)}:=\frac1{d_j}\mathbf 1^\top Z^{(j)}\), then
\begin{equation}\label{eq:center-proj}
\Big(I_{d_j}-\frac1{d_j}\mathbf 1\mathbf 1^\top\Big)Z^{(j)}
=Z^{(j)}-\bar Z^{(j)}\mathbf 1
\sim \mathcal N\Big(0,\ I_{d_j}-\frac1{d_j}\mathbf 1\mathbf 1^\top\Big),\nonumber
\end{equation}
so
\begin{equation}\label{eq:max-shift}
\max_{k\le d_j} (h^{(j)})_k\ \stackrel{d}{=}\ \max_{k\le d_j}\big((Z^{(j)})_k-\bar Z^{(j)}\big)
=\max_{k\le d_j} (Z^{(j)})_k-\bar Z^{(j)}.\nonumber
\end{equation}
Taking expectations and using \(\mathbb E[\bar Z^{(j)}]=0\),
\begin{equation}\label{eq:max-equality}
\mathbb E\Big[\max_{k\le d_j} (h^{(j)})_k\Big]
=
\mathbb E\Big[\max_{k\le d_j} Z_k\Big]\leq \mathcal{O}(\sqrt{\log d_i}),
\qquad Z_k\stackrel{iid}{\sim}\mathcal N(0,1).\nonumber
\end{equation}
Thus, due to Eq. \eqref{eq:width-reduce} and the calculations above, we have
\begin{equation}\label{eq:width-final}
w(\mathcal{X})\leq O\left(\sum_{j=1}^m \sqrt{d_j\log d_j}\right).\nonumber
\end{equation}

\section{Lower Bounds for Structured Sets}\label{appendix-structured-sets-lower-bound}
First, in section \ref{appendix-multi-task-lower bound}, we present the lower bound for the multi-task MAB problem. Next, in section \ref{appendix-multi-task-lower bound}, we present the lower bound for hypercubes. Finally, in section \ref{appendix-m-sets-lower-bound}, we present the lower bound for $m$-sets.
\subsection{Lower Bound for Multi-task MAB}\label{appendix-multi-task-lower bound}

%\subsection{Lower Bound for Multi-task MAB}


%\textbf{Formal proof}:
Recall that in the Multi-task MAB, the set of arms $\mathcal{X}$ is defined as follows: 
\begin{equation*}
    \mathcal{X}=\left\{\mathbf{x}\in \{0,1\}^d: \forall j\in [m] \sum_{i=d_{1:j-1}+1}^{d_{1:j}} \mathbf{x}[i]=1\right\}
\end{equation*}
where $d_i\geq 2$, $d=\sum_{i=1}^md_i$, $d_{1:j}=\sum_{i=1}^jd_i$ and $d_{1:0}=0$ for all $j\in [m]$.

We first focus on the case when $d_j\geq 100$ for all $j\in[m]$.

Let us fix one regret minimizing algorithm, say $\mathcal{A}$ and assume that $\mathcal{A}$ is deterministic given the reward realizations. Let us assume that $\mathcal{A}$ terminates after taking $T:=\frac{(\sum_{i=1}^m\sqrt{d_i})^2}{20000\varepsilon^2}$ samples (if the algorithm terminates before this, we can take additional dummy samples). Let $\widehat{x}\in \mathcal{X}$ be the arm recommended by $\mathcal{A}$ after sampling the arms $\mathbf{x}_1,\ldots,\mathbf{x}_T$ and terminating. 

Later in this section, we construct a set of input instances $\mathcal{I}$. Each instance $I\in \mathcal{I}$ is associated with a reward vector $\theta_I\in \mathbb{R}^d$ such that the following hold:
\begin{itemize}
    \item If $\mathcal{A}$ samples an arm $x\in \mathcal{X}$, it observes $\langle x,\theta_I\rangle+\eta$ where $\eta\sim \mathcal{N}(0,1)$.
    \item For all $x\in \mathcal{X}$, we have $\langle x,\theta_I\rangle\geq 0$.
    \item We have $\max_{x\in \mathcal{X}}\langle x,\theta_I\rangle=10\varepsilon$.
\end{itemize}

Suppose we show that $\mathbb{E}_{I\sim Unif(\mathcal{I})}[\langle \widehat{x}, \theta_I\rangle]\leq 7\varepsilon$, where $\mathbb E_{I}[\cdot]$ is the expectation under the instance $I$. Then by Yao's lemma, we have for any randomized algorithm, the recommended arm $\widehat{x}$ satisfies $\min_{I\in \mathcal{I}}\mathbb E_{I}[\langle \widehat{x},\theta_I\rangle]\leq 7\varepsilon$. Then due to markov's inequality, we have that with probability at least $0.1$, there exists an instance on which the randomized algorithm does not recommend an $\varepsilon$-best arm.

Now we focus on showing that $\mathbb{E}_{I\sim Unif(\mathcal{I})}[\langle \widehat{x}, \theta_I\rangle]\leq 7\varepsilon$. We begin with construction instances $I$ with their corresponding reward vector $\theta_I$. These instances include both the set of all input instances and a set of alternate instances that will be used to argue the performance of the algorithm $\mathcal{A}$.

 % Now we show that algorithm $\mathcal{A}$ incurs a regret of $\Omega(\sum_{i=1}^m\sqrt{d_iT})$. We later extend the result to randomized algorithms using Yao's lemma. For all $j\in[m]$, let $\varepsilon_j>0$ be a parameter that we fix later in the proof.

Let $\tilde{\mathcal{X}}=\left\{\mathbf{x}\in \{0,1\}^d: \forall j\in [m] \sum_{i=d_{1:j-1}+1}^{d_{1:j}} \mathbf{x}[i]\leq 1\right\}$.
First we describe an instance $I_{\tilde{\mathbf{x}}}$, where $\tilde{\mathbf{x}}\in \tilde{\mathcal{X}}$. In this instance, we define the associated reward vector $\theta_{I_{\tilde{\mathbf{x}}}}$ as follows. For all $j\in [m]$ and $i\in \{d_{1:j-1}+1,\ldots,d_{1:j}\}$, we have $\theta_{I_{\tilde{\mathbf{x}}}}[i]=10\varepsilon_j\cdot \tilde{\mathbf{x}}[i]$ where $\varepsilon_j=\frac{\varepsilon\cdot\sqrt{d_j}}{\sum_{s\in[m]}\sqrt{d_s}}$. Observe that for all $x\in\mathcal{X}$, we have $\langle x,\theta_{I_{\tilde x}}\rangle\geq 0$. Also observe that $\max_{x\in\mathcal{X}}\langle x,\theta_{I_{\tilde x}}\rangle=10\varepsilon$ if $\tilde{x}\in \mathcal{X}$.

Let $r^{(j)}_{\tilde x}:=\mathbb{E}_{I_{\tilde x}}[\sum_{i=d_{1:j-1}+1}^{d_{1:j}}\widehat{x}[i]\cdot \theta_{I_{\tilde x}}[i]]$. Observe that $\mathbb{E}_{I_{\tilde x}}[\langle \widehat{x},\theta_{I_{\tilde x}}\rangle]=\sum_{j=1}^m r^{(j)}_{\tilde x}$.

%If for all  $\mathbb{E}_{I_x}[\sum_{i=d_{1:j-1}+1}^{d_{1:j}}\widehat{x}[i]\cdot \theta_{I_{\tilde x}}[i]]$.
    
    
    Let $\mathcal{I}=\bigcup_{\mathbf{x}\in\mathcal{X}}I_{\mathbf{x}}$. Fix an index $j\in[m]$. We now show that $\mathbb{E}_{I_{\mathbf{x}'}\sim \mathrm{Unif}(\mathcal{I})}[\sum_{i=d_{1:j-1}+1}^{d_{1:j}}\widehat{x}[i]\cdot \theta_{I_{\mathbf{x}'}}[i]]\leq 7\varepsilon_j$. Let $$\mathcal{X}^{(j)}:=\Bigg\{\mathbf{x}\in \{0,1\}^d: \forall i\in[m]\setminus\{j\}\; \sum_{s=d_{1:i-1}+1}^{d_{1:i}}\mathbf{x}[s]=1,\; \sum_{s=d_{1:j-1}+1}^{d_{1:j}}\mathbf{x}[s]=0\Bigg\}.$$ For any $\mathbf{x}\in \mathcal{X}^{(j)}$, let $\mathbf{x}^{(i)}$ be the vector in $\mathcal{X}$ such that $\mathbf{x}^{(i)}[s]=\mathbf{x}[s]$ for all $s\notin \{d_{1:j-1}+1,\ldots,d_{1:j}\}$ and $\mathbf{x}^{(i)}[d_{1:j-1}+i]=1$.

    Let us fix $\mathbf{x}\in \mathcal{X}^{(j)}$. Now we claim that there is a set $\mathcal{S}_{\mathbf{x}}\subseteq[d_j]$ with at least $d_j/3$ indices such that for each $i\in \mathcal{S}_\mathbf{x}$, we have $r^{(j)}_{\mathbf{x}^{(i)}} \leq \varepsilon_j$.
    
    Before we prove our claim, we first show that if our claim holds true, then we have that $$\mathbb{E}_{I_{\mathbf{x}'}\sim \mathrm{Unif}(\mathcal{I})}[\sum_{i=d_{1:j-1}+1}^{d_{1:j}}\widehat{x}[i]\cdot \theta_{I_{\mathbf{x}'}}[i]]\leq 7\varepsilon_j.$$ Now we have the following:
    \begin{align*}
        \mathbb{E}_{I_{\mathbf{x}'}\sim \mathrm{Unif}(\mathcal{I})}[\sum_{i=d_{1:j-1}+1}^{d_{1:j}}\widehat{x}[i]\cdot \theta_{I_{\mathbf{x}'}}[i]]&=\frac{1}{\prod_{s=1}^md_s}\sum_{\mathbf{x}\in\mathcal{X}^{(j)}}\sum_{i=1}^{d_j} r^{(j)}_{\mathbf{x}^{(i)}} \\
        &= \frac{1}{\prod_{s=1}^md_s}\sum_{\mathbf{x}\in\mathcal{X}^{(j)}}\left(\sum_{i\in \mathcal{S}_{\mathbf{x}}}r^{(j)}_{\mathbf{x}^{(i)}}+\sum_{i\in [d_j]\setminus\mathcal{S}_{\mathbf{x}}}r^{(j)}_{\mathbf{x}^{(i)}}\right)\\
        &\leq \frac{1}{\prod_{s=1}^md_s}\sum_{\mathbf{x}\in\mathcal{X}^{(j)}}\left(\frac{d_j}{3}\cdot \varepsilon_j+\frac{2d_j}{3}\cdot 10\varepsilon_j\right)\\
        &\leq \frac{1}{\prod_{s=1}^md_s}\sum_{\mathbf{x}\in\mathcal{X}^{(j)}}7d_j\varepsilon_j  \\
        & = \frac{7\varepsilon_j}{\prod_{s=1}^md_s} \cdot \prod_{s\neq j}d_s\cdot d_j \\
        & = 7\varepsilon_j
    \end{align*}

    Now we prove our claim. We use the adaptive KL chain rule (Lemma \ref{kl-chain-rule}) in our analysis.
    

    Let $\mathbb{P}_I$ denote the probability law instances under an instance $I$. For an instance $I_{\mathbf{x}^{(i)}}$, let $f_i(\ell_1,\ldots,\ell_T)$ denote the joint PDF for the tuple of reward values observed by $\mathcal{A}$ in each round under the probability law $\mathbb{P}_{I_{\mathbf{x}^{(i)}}}$. Observe that our sample space is $\Omega=\mathbb{R}^T$. This is a valid sample space as $\mathcal{A}$ is deterministic and the probability density function of the reward values in round $t$ only depends on the reward values it observed in the previous rounds. Similarly for the alternate instance  $I_{\mathbf{x}}$, let $f_0(\ell_1,\ldots,\ell_T)$ denote the joint PDF for the tuple of loss values observed by $\mathcal{A}$ in each round under the probability law $\mathbb{P}_{I_{\mathbf{x}}}$.

% \begin{figure}[h]
%     \centering
%     \includegraphics[width=\textwidth]{Images-3.pdf} % or output.eps
%     %\caption{Instance $I_S$ where $S=(j_1,j_2,\ldots,j_K)$}
%     %\label{fig:yourlabel}
%     \end{figure}
    First observe that the instances $I_{\mathbf{x}^{(i)}}$ and $I_{\mathbf{x}}$ only differ at index $d_{1:j-1}+i$. For each $\omega\in \Omega$, let $\mathbf{x}_{1,\omega},\mathbf{x}_{2,\omega},\ldots, \mathbf{x}_{T,\omega}$ be the sequence of arms chosen by $\mathcal{A}$ on $\omega$.
    Conditioning on a set of outcomes $X_1=\omega_1,X_2=\omega_2,\ldots,X_{t-1}=\omega_{t-1}$, we have $X_t\sim \mathcal{N}(\mu_i,1)$ for the instance $I_{\mathbf{x}^{(i)}}$ and $X_t\sim \mathcal{N}(\mu_0,1)$ for the instance $I_{\mathbf{x}}$ where $\mu_i-\mu_0=10\varepsilon_j\cdot \mathbf{x}_{t,\omega}[d_{1:j-1}+i]$. Let $T_i=\sum_{t=1}^T\mathbf{x}_t[d_{1:j-1}+i]$.  For each $\omega\in \Omega$, let $T_{i,\omega}=\sum_{t=1}^T\mathbf{x}_{t,\omega}[d_{1:j-1}+i]$. Note that $T_i$ is a random variable and $T_{i,\omega}$ is a fixed value. Let $X_{-t}$ denote $X_1,\ldots,X_{t-1}$ and $\omega_{-t}$ denote $\omega_1,\ldots,\omega_{t-1}$. Now we have the following:
    \begin{align*}
         KL(f_0,f_i)&=\int\limits_{\omega\in \Omega}f_0(\omega)\left(KL(f_0(X_1),f_i(X_1))+\sum_{t=2}^T KL(f_0(X_t|X_{-t}=\omega_{-t}),f_i(X_t|X_{-t}=\omega_{-t}))\right)d\omega\\
         &=50\varepsilon_j^2\int\limits_{\omega\in \Omega}f_0(\omega)\sum_{t=1}^T\mathbf{x}_{t,\omega}[d_{1:j-1}+i]\;d\omega\\
         &=50\varepsilon_j^2 \int\limits_{\omega\in \Omega}f_0(\omega)T_{i,\omega}\;d\omega\\
         &= 50\varepsilon_j^2\cdot \mathbb{E}_{I_{\mathbf{x}}}[T_{i}]
    \end{align*}

    Now observe that $\sum_{i=1}^{d_j}\mathbb{E}_{I_{\mathbf{x}}}[T_i]=T$. Hence, there exists a set $\mathcal{S}_\mathbf{x}^{(1)}\subseteq[d_j]$ with at least $2d_j/3$ indices such that for each $i\in \mathcal{S}_x^{(1)}$, we have $\mathbb{E}_{I_{\mathbf{x}}}[T_i]\leq \frac{3T}{d_j}$. Similarly, there exists a set $\mathcal{S}_\mathbf{x}^{(2)}\subseteq[d_j]$ with at least $2d_j/3$ indices such that for each $i\in \mathcal{S}_x^{(2)}$, we have $\mathbb{P}_{I_{\mathbf{x}}}[\widehat{x}[d_{1:j-1}+i]=1]\leq \frac{3}{d_j}$. Now we define $\mathcal{S}_\mathbf{x}:=\mathcal{S}_x^{(1)}\cap \mathcal{S}_x^{(2)}$. Observe that $|\mathcal{S}_\mathbf{x}|\geq d_j/3$. Next for each $i\in \mathcal{S}_x$, we have $\mathbb{E}_{I_{\mathbf{x}}}[T_i]\leq \frac{3T}{d_j}$ and $\mathbb{P}_{I_{\mathbf{x}}}[\widehat{x}[d_{1:j-1}+i]=1]\leq \frac{3}{d_j}$. Now for each $i\in \mathcal{S}_\mathbf{x}$, we have $KL(f_0,f_i)\leq\frac{150\varepsilon_j^2T}{d_j}=\frac{3}{400}$. 

    Fix $i\in \mathcal{S}_x$. Let $A_i$ be the event that $\widehat{x}[d_{1:j-1}+i]=1$. Now due to Pinsker's inequality we have the following:
    \begin{align*}
        \mathbb{P}_{I_{\mathbf{x}^{(i)}}}(A_i)&\leq  \mathbb{P}_{I_{\mathbf{x}}}(A_i)+\sqrt{\frac{KL(f_0,f_i)}{2}}\\
        &\leq \frac{3}{d_j}+\sqrt{\frac{3}{800}}\\
        &\leq\frac{3}{100}+\sqrt{\frac{3}{800}}\tag{as $d_j\geq 100$}\\
        &< \frac{1}{10}
    \end{align*}

    Hence, we have $r^{(j)}_{\mathbf{x}^{(i)}} < \frac{1}{10}\cdot 10\varepsilon_j= \varepsilon_j$.

    Now, we conclude the proof by showing that $\mathbb{E}_{I\sim Unif(\mathcal{I})}[\langle \widehat{x}, \theta_I\rangle]\leq 7\varepsilon$:
    \begin{align*}
        \mathbb{E}_{I\sim Unif(\mathcal{I})}[\langle \widehat{x}, \theta_I\rangle]&=\sum_{j=1}^m \mathbb{E}_{I_{\mathbf{x}'}\sim \mathrm{Unif}(\mathcal{I})}[\sum_{i=d_{1:j-1}+1}^{d_{1:j}}\widehat{x}[i]\cdot \theta_{I_{\mathbf{x}'}}[i]]\\
        &=7\sum_{j=1}^m \varepsilon_j\\
        &=7\varepsilon\cdot\frac{\sum_{j=1}^m \sqrt{d_j}}{\sum_{s=1}^m \sqrt{d_s}}\tag{as $\varepsilon_j=\frac{\varepsilon\cdot\sqrt{d_j}}{\sum_{s=1}^m \sqrt{d_s}}$}\\
        &=7\varepsilon
    \end{align*}

\subsubsection{$d_j<100$ Case}

For the simplicity of presentation let us focus on the case $d_j=2$. This analysis can be easily extended to the other cases easily.
For the case $d_j=2$, the construction of the hard instances remain exactly the same as that of the $d_j\geq 100$ case. The analysis for this case only differs at the end. Consider an index $i\in\{1,2\}$. Observe that $KL(f_0,f_i)\leq50\varepsilon_j^2T=\frac{1}{200}$. Recall that $A_i$ is the event that $\widehat{x}[d_{1:j-1}+i]=1$. Now due to Pinsker's inequality we have the following:
    \begin{align*}
        \mathbb{P}_{I_{\mathbf{x}^{(i)}}}(A_i)\leq  \mathbb{P}_{I_{\mathbf{x}}}(A_i)+\sqrt{\frac{KL(f_0,f_i)}{2}}\leq \mathbb{P}_{I_{\mathbf{x}}}(A_i)+\frac{1}{20}\\
    \end{align*}
Recall that $r^{(j)}_{\tilde x}:=\mathbb{E}_{I_{\tilde x}}[\sum_{i=d_{1:j-1}+1}^{d_{1:j}}\widehat{x}[i]\cdot \theta_{I_{\tilde x}}[i]]$. Now we have the following:
    \begin{align*}
        \mathbb{E}_{I_{\mathbf{x}'}\sim \mathrm{Unif}(\mathcal{I})}[\sum_{i=d_{1:j-1}+1}^{d_{1:j}}\widehat{x}[i]\cdot \theta_{I_{\mathbf{x}'}}[i]]&=\frac{1}{\prod_{s=1}^md_s}\sum_{\mathbf{x}\in\mathcal{X}^{(j)}}\sum_{i=1}^{2} r^{(j)}_{\mathbf{x}^{(i)}} \\
        &\leq \frac{10\varepsilon_j}{\prod_{s=1}^md_s}\sum_{\mathbf{x}\in\mathcal{X}^{(j)}}\left(\mathbb{P}_{I_{\mathbf{x}}}(A_1)+\mathbb{P}_{I_{\mathbf{x}}}(A_2)+0.1\right)\\
        & = \frac{10\varepsilon_j}{\prod_{s=1}^md_s} \cdot\frac{1.1}{2}\cdot \prod_{s\neq j}d_s\cdot d_j \\
        & = 5.5\varepsilon_j
    \end{align*}

\subsection{Hypercubes}\label{appendix-hypercubes-lower-bound}
For the $\{0,1\}^d$ hypercube, we consider the family of hard instances $\{I_{\mathbf{x}}\}_{\mathbf{x}\in\{0,1\}^d}$ with corresponding reward vectors $\{\theta_{\mathbf{x}}\}_{\mathbf{x}\in\{0,1\}^d}$. For each $\mathbf{x}\in\{0,1\}^d$, define $\theta_{\mathbf{x}}\in\mathbb{R}^d$ coordinate-wise by
$(\theta_{\mathbf{x}})_i=\mathbbm{1}\{\mathbf{x}_i=1\}\cdot \frac{10\varepsilon}{d}-\mathbbm{1}\{\mathbf{x}_i=0\}\cdot \frac{10\varepsilon}{d}$.
An analysis analogous to our $d_j=2$ argument for multi-task MAB yields a lower bound of $\Omega(d^2/\varepsilon^2)$. The only minor difference is that, in the multi-task setting, we reasoned about the expected value of $\langle x,\theta\rangle$, whereas here one must instead work with the expected gap $\langle x,\theta\rangle-\min_{x\in\{0,1\}^d}\langle x,\theta\rangle$.

For the $\{-1,+1\}^d$ hypercube, we similarly consider the family of hard instances $\{I_{\mathbf{x}}\}_{\mathbf{x}\in\{-1,+1\}^d}$ with corresponding reward vectors $\{\theta_{\mathbf{x}}\}_{\mathbf{x}\in\{-1,+1\}^d}$. For each $\mathbf{x}\in\{-1,+1\}^d$, define $\theta_{\mathbf{x}}\in\mathbb{R}^d$ coordinate-wise by
$(\theta_{\mathbf{x}})_i=\mathbf{x}_i\cdot \frac{5\varepsilon}{d}$.
An analysis analogous to our $d_j=2$ argument for multi-task MAB again yields a lower bound of $\Omega(d^2/\varepsilon^2)$. As above, the only difference is that the argument proceeds via the expected gap $\langle x,\theta\rangle-\min_{x\in\{-1,+1\}^d}\langle x,\theta\rangle$ rather than expectation of $\langle x,\theta\rangle$ itself.




\subsection{$m$-Sets Lower Bound}\label{appendix-m-sets-lower-bound}
\begin{theorem}\label{appendix:thm-m-sets}
Let us denote the $m$-sets by
\[
\mathcal X=\{x\in\{0,1\}^d:\|x\|_0=m\}.
\]
In each round $t=1,\dots,T$, the learner chooses $x_t\in\mathcal X$ and observes
\[
y_t=\langle x_t,\theta\rangle+\eta_t,\qquad \eta_t\sim \mathcal N(0,1)\ \text{i.i.d.}
\]
Suppose $d-m+1\ge 20m$. Then there exists a universal constant $c>0$ such that for any (possibly adaptive and randomized) algorithm,
if
\[
T \le c\,\frac{m(d-m+1)}{\varepsilon^2},
\]
then there exists an instance $\theta$ for which the algorithm fails to output an $\varepsilon$-optimal arm with probability
at least $2/9$.
\end{theorem}

\begin{proof}
We first describe the family of hard instances. For each $S\subset[d]$ with $|S|=m$, define
\[
\theta^{(S)}_i=
\begin{cases}
\Delta & i\in S,\\
0 & i\notin S,
\end{cases}
\qquad\text{where}\qquad \Delta:=\frac{10\varepsilon}{m}.
\]
Let $S$ be uniform over all $m$-sized subsets of $[d]$ and the environment parameter be $\theta=\theta^{(S)}$.
For any output $\hat x\in\mathcal X$, write $\hat S:=\mathrm{supp}(\hat x)$ so $|\hat S|=m$. Then
\[
\langle \hat x,\theta^{(S)}\rangle=\Delta\,|\hat S\cap S|.
\]
The optimal value is $\max_{x\in\mathcal X}\langle x,\theta^{(S)}\rangle=\Delta|S|=m\Delta=10\varepsilon$.
Thus on any instance, $\varepsilon$-success implies $\langle \hat x,\theta^{(S)}\rangle\ge 9\varepsilon$.

By Yao's lemma it suffices to fix an arbitrary deterministic algorithm and analyze its error under the random draw of $S$.

Given $S$, let $I$ be uniform on $S$ (an auxiliary random variable), and define $B:=S\setminus\{I\}$ so $|B|=m-1$.
Conditioning on $(S,\hat S)$,
\[
\mathbb P(I\in \hat S\mid S,\hat S)=\frac{|\hat S\cap S|}{|S|}=\frac{|\hat S\cap S|}{m},
\]
so by the tower rule
\begin{equation}\label{eq:overlap-identity}
\mathbb E[|\hat S\cap S|]=m\,\mathbb P(I\in \hat S).
\end{equation}
Therefore,
\begin{equation}\label{eq:value-vs-I}
\mathbb E[\langle \hat x,\theta^{(S)}\rangle]
=\Delta\,\mathbb E[|\hat S\cap S|]
=m\Delta\,\mathbb P(I\in \hat S)
=10\varepsilon\cdot \mathbb P(I\in \hat S).
\end{equation}
So it is enough to show $\mathbb P(I\in \hat S)\le 0.7$, which would imply
$\mathbb E[\langle \hat x,\theta^{(S)}\rangle]\le 7\varepsilon$.

% \medskip
% \textbf{A coupling lemma (how to condition on $B$).}
% Let $n:=d-m+1$. Note that $|[d]\setminus B|=n$.

We now prove the following technical lemma.
\begin{lemma}[Equivalent sampling]\label{appendix-lem-equivalent-sampling}
The joint law of $(B,I,S)$ defined above is the same as:
(i) draw $B$ uniformly among all $(m-1)$-subsets of $[d]$,
(ii) draw $I$ uniformly from $[d]\setminus B$,
(iii) set $S=B\cup\{I\}$.
In particular, conditional on $B$, $I\sim \mathrm{Unif}([d]\setminus B)$.
\end{lemma}
\begin{proof}
Fix $b\subset[d]$ with $|b|=m-1$ and $i\notin b$.
Under the original procedure,
\[
\mathbb P(B=b,I=i)
=\mathbb P(S=b\cup\{i\})\cdot \mathbb P(I=i\mid S=b\cup\{i\})
=\frac{1}{\binom{d}{m}}\cdot \frac{1}{m}.
\]
Under the alternative procedure,
\[
\mathbb P(B=b,I=i)
=\mathbb P(B=b)\cdot \mathbb P(I=i|B=b)=\frac{1}{\binom{d}{m-1}}\cdot \frac{1}{d-m+1}=\frac{1}{\binom{d}{m}}\cdot \frac{1}{m}.
\]
Hence, both sampling procedures have the same joint law.
\end{proof}

Fix any $B$ and define the alternative instance $\theta^{(B)}$ by
$\theta^{(B)}_j=\Delta$ if $j\in B$ and $0$ otherwise.
For each $i\notin B$, define
\[
N_i:=\sum_{t=1}^T x_t(i)\qquad\text{and}\qquad A_i:=\{i\in \hat S\}.
\]
Because every played action has exactly $m$ ones,
\[
\sum_{i=1}^d N_i=mT \quad\Rightarrow\quad \sum_{i\notin B}\mathbb E_{\theta^{(B)}}[N_i]\le mT.
\]
Because $|\hat S|=m$ always,
\[
\sum_{i\notin B}\mathbb P_{\theta^{(B)}}(A_i)\le m.
\]
Let $n:=d-m+1=|[d]\setminus B|$. Due to the above two inequalities, there exists a set
$G_B\subseteq [d]\setminus B$ with $|G_B|\ge n/2$ such that for all $i\in G_B$,
\begin{equation}\label{eq:good-bounds}
\mathbb E_{\theta^{(B)}}[N_i]\le \frac{4mT}{n},
\qquad
\mathbb P_{\theta^{(B)}}(A_i)\le \frac{4m}{n}.
\end{equation}

Fix $i\in G_B$. Compare $\theta^{(B)}$ with $\theta^{(B\cup\{i\})}$.
They differ only in coordinate $i$ by $\Delta$.

Let $P_B$ and $P_{B\cup\{i\}}$ be the probability law under
$\theta^{(B)}$ and $\theta^{(B\cup\{i\})}$, respectively.
For Gaussian noise with variance $1$, the one-step KL between $\mathcal N(\mu,1)$ and $\mathcal N(\mu+\Delta,1)$
equals $\Delta^2/2$.
Using the adaptive KL chain rule (Lemma \ref{kl-chain-rule}) in the same way as we did the multi-task MAB, we get
\begin{equation}\label{eq:kl}
\mathrm{KL}(P_B\|P_{B\cup\{i\}})
=\frac{\Delta^2}{2}\,\mathbb E_{\theta^{(B)}}[N_i]
\le \frac{\Delta^2}{2}\cdot \frac{4mT}{n}
=2\Delta^2\frac{mT}{n},
\end{equation}
where we used \eqref{eq:good-bounds}.
Pinsker's inequality gives, for the event $A_i$,
\[
\mathbb P_{\theta^{(B\cup\{i\})}}(A_i)
\le \mathbb P_{\theta^{(B)}}(A_i) + \sqrt{\tfrac12\,\mathrm{KL}(P_B\|P_{B\cup\{i\}})}.
\]
Combining with \eqref{eq:good-bounds} and \eqref{eq:kl},
\begin{equation}\label{eq:pinsker-bound}
\mathbb P_{\theta^{(B\cup\{i\})}}(i\in \hat S)
\le \frac{4m}{n} + \Delta\sqrt{\frac{mT}{n}}.
\end{equation}
Assume $n\ge 20m$ so that $4m/n\le 0.2$.
Also assume  $T \le \frac{1}{2500}\cdot \frac{mn}{\varepsilon^2}.$ Since $\Delta=10\varepsilon/m$, we have
\[
\Delta\sqrt{\frac{mT}{n}}
=\frac{10\varepsilon}{m}\sqrt{\frac{mT}{n}}
\le \frac{10\varepsilon}{m}\sqrt{\frac{m}{n}\cdot \frac{1}{2500}\frac{mn}{\varepsilon^2}}
=0.2.
\]
Plugging into \eqref{eq:pinsker-bound} yields, for every $i\in G_B$,
\begin{equation}\label{eq:good-i-prob}
\mathbb P_{\theta^{(B\cup\{i\})}}(i\in \hat S)\le 0.4.
\end{equation}

Due to Lemma \ref{appendix-lem-equivalent-sampling}, conditional on $B$, $I$ is uniform on $[d]\setminus B$, and if $I=i$ then the instance is
$\theta^{(B\cup\{i\})}$. Therefore, we have:
\begin{equation}\label{eq:avg}
\mathbb P(I\in \hat S\mid B)
=\frac{1}{n}\sum_{i\in[d]\setminus B}\mathbb P_{\theta^{(B\cup\{i\})}}(i\in \hat S).
\end{equation}
Split the sum into $G_B$ and its complement. As $|G_B|\ge n/2$,
applying \eqref{eq:good-i-prob} for each $i\in G_B$, and using the trivial upper bound of $1$ for $i\notin G_B$, we get:
\[
\mathbb P(I\in \hat S\mid B)
\le \frac{1}{n}\Big(\frac{n}{2}\cdot 0.4 + \frac{n}{2}\cdot 1\Big)=0.7.
\]
Averaging over $B$ gives
\begin{equation}\label{eq:PI}
\mathbb P(I\in \hat S)\le 0.7.
\end{equation}

From \eqref{eq:value-vs-I} and \eqref{eq:PI}, we get:
\[
\mathbb E[\langle \hat x,\theta^{(S)}\rangle]\le 10\varepsilon\cdot 0.7=7\varepsilon.
\]
On the other hand, success implies $\langle \hat x,\theta^{(S)}\rangle\ge 9\varepsilon$, and the value is always
nonnegative, hence
\[
\mathbb E[\langle \hat x,\theta^{(S)}\rangle]\ge 9\varepsilon\cdot \mathbb P(\text{success}).
\]
Therefore $\mathbb P(\text{success})\le 7/9$, so $\mathbb P(\text{fail})\ge 2/9$ under uniform distribution over the hard instances.
By Yao's principle, there exists a fixed instance $\theta$ on which the algorithm fails with probability at least $2/9$.
This proves the theorem with a suitable universal constant $c$ (for instance $c=1/2500$).
\end{proof}
\section{Unit Ball Lower Bound}\label{appendix-ball-lower-bound}
Fix $\varepsilon\in (0,1]$, $\delta\in(0,1)$ and an $(\varepsilon,\delta)$-PAC algorithm. Let the algorithm run for $T:=\frac{H_1\log(1/\delta)+H_2}{\varepsilon^2}$ rounds and output $\hat x\in \mathcal{X}:=\mathbb{B}_d$.
Define the simple regret for any $\theta$ as
\[
R_T(\theta)\;:=\;\max_{x\in \mathcal{X}}\langle x,\theta\rangle-\langle \hat x,\theta\rangle \;\ge 0.
\]
Observe that $R_T(\theta)\leq 2\|\theta\|$.


As the algorithm is $(\varepsilon,\delta)$-PAC, we have $\mathbb P(R_T(\theta)>\varepsilon)\le \delta$. Hence, taking expectation, we have
\begin{equation}\label{eq:ball-simple-regret-upper}
    \mathbb{E}[R_T(\theta)]\leq \varepsilon+2\|\theta\|\cdot\delta 
\end{equation}


Consider a small enough $\varepsilon$ such that $T\geq d^2$.
Now due to \cite{chen2024assouad}, there exists a $\theta\in\mathbb{R}^d$ such that $\|\theta\|_2\leq c_0d/\sqrt{T}$ and $\mathbb E[R_T(\theta)]\geq c_1d/\sqrt{T}$ for some absolute constants $c_0,c_1\in (0,1)$. Let us consider now consider such a $\theta$. If we choose $\delta=c_1/4$, we have $\varepsilon\geq \frac{c_1d}{2\sqrt{T}}$ due to \eqref{eq:ball-simple-regret-upper}. As $H_1\leq H_2$ and $T=\frac{H_1\log(1/\delta)+H_2}{\varepsilon^2}$, we have $H_2\geq c_2 d^2$ for some absolute constant $c_2>0$.


Let us also consider the case where stopping time $\tau$ is randomized as it will be used in the next section. Due to \cite{chen2024assouad}, there exists a $\theta\in\mathbb{R}^d$ such that $\|\theta\|_2\leq c_0d/\sqrt{T}$ and $\mathbb E[R_T(\theta)]\geq c_1d/\sqrt{T}$ for some absolute constants $c_0,c_1\in (0,1)$. Consider one such $\theta$. Let assume that $\mathbb{E}[\tau]=T_0:=\frac{H_1\log(1/\delta)+H_2}{\varepsilon^2}$ with $0<H_1\leq H_2$. Let us consider $T=\frac{1}{\delta}\cdot T_0 $. If the algorithm terminates before $T$ rounds, we pad the remaining rounds using dummy samples. Now we have the following due to Markov's inequality:
\[
   \mathbb E[R_T(\theta)]\leq \varepsilon +2 \|\theta\|( \mathbb{P}(\tau>T)+\mathbb{P}(R_T(\theta)>\varepsilon,\tau\leq T))\leq \varepsilon +4\delta\cdot \|\theta\|
\]
Hence, for small enough absolute constant $\delta$, we have $H_2\geq c_3\cdot d^2$ for an absolute constant $c_3$.

%\section{Polynomial separation instance's non-adaptive lower bound}
% Recall that $\mathcal{X}:=\bigcup_{i=1}^k \mathcal{X}_i$ where each $\mathcal{X}_i$ is a $d$-dimensional $\ell_2$-ball. Consider an optimal non-adaptive $(\varepsilon,\delta)$-PAC algorithm that takes $T=\frac{H_1\log(1/\delta)+H_2}{\varepsilon^2}$ samples with $0<H_1\leq H_2$. Due to our non-adaptive upper bound algorithm and Theorem~\ref{thm:adaptive-lower-bound}, we have $H_1=\Theta(d)$. 

% Now let us assume that $T\leq c\cdot \left(\frac{kd\log(1/\delta)+kd^2}{\varepsilon^2}\right)$. Let $\tau_i$ be a random variable that denotes the number of samples allocated by the algorithm to $\mathcal{X}_i$. There exists $i\in [k]$ such $T_i:=\mathbb{E}[\tau_i]\leq c\cdot \left(\frac{d\log(1/\delta)+d^2}{\varepsilon^2}\right)$. Consider a hard instance where every other block is assigned zero vector and the block containing $\mathcal{X}_i$ is the assigned the hard instance from the previous section. Now if the algorithm outputs a vector $\hat x$ such that it belongs to a block not containing $\mathcal{X}_i$, then it is equivalent to choosing the zero vector from $\mathcal{X}_i$. Hence, the lower bound result from the previous section is applicable on $\mathcal{X}_i$. However, if the constant $c$ is sufficiently small, then we would be contradicting the result from the previous section. Hence, $H_2\geq c'\cdot kd^2$ for some absolute constant $c'>0$. 

\section{Polynomial Separation Instance's Non-Adaptive Lower Bound}\label{appendix-poly-sep-lower-bound}
Recall $\mathcal{X}=\bigcup_{i=1}^k \mathcal{X}_i\subseteq \mathbb{R}^{kd}$, where for each $i\in[k]$ we define the \emph{$i$-th block of coordinates} as $B_i:=\{(i-1)d+1,\ldots,id\}$ and
\[
\mathcal{X}_i:=\Bigl\{x\in\mathbb{R}^{kd}:\ \mathrm{supp}(x)\subseteq B_i,\ \|x\|_2\le 1\Bigr\}.
\]
Consider any (possibly randomized) non-adaptive $(\varepsilon,\delta)$-PAC algorithm that uses $T=\frac{H_1\log(1/\delta)+H_2}{\varepsilon^2}$ samples with $0<H_1\le H_2$, and let $\tau_i:=\sum_{t=1}^T \mathbbm{1}\{X_t\in\mathcal{X}_i\}$ denote the (random) number of samples allocated to block $i$, so that $\sum_{i=1}^k \tau_i=T$ and hence there exists $i^\star\in[k]$ with $T_{i^\star}:=\mathbb{E}[\tau_{i^\star}]\le T/k$. Fix such an $i^\star$ and define a hard instance $\theta\in\mathbb{R}^{kd}$ supported only on block $B_{i^\star}$ as follows: set $\theta_j=0$ for all $j\notin B_{i^\star}$, and on the coordinates in $B_{i^\star}$ set $(\theta_j)_{j\in B_{i^\star}}=\vartheta$ where $\vartheta\in\mathbb{R}^d$ is chosen according to the unit-ball lower bound from the previous section (\cite{chen2024assouad} together with the padding and Markov's inequality argument), so that for some absolute constants $c_0,c_1\in(0,1)$ one has $\|\vartheta\|_2\le c_0\cdot \sqrt{\delta}\cdot d/\sqrt{T_{i^\star}}$ and any algorithm that receives $T_{i^\star}/\delta$ informative samples on this block satisfies 
\[
c_1\cdot \sqrt{\delta}\cdot d/\sqrt{T_{i^\star}} \le\mathbb{E}[R_{T_{i^\star}/\delta}(\vartheta)]\le \varepsilon +4\delta\cdot\|\vartheta\| 
\]
For our chosen $\theta$, any action $x\in\mathcal{X}_j$ with $j\neq i^\star$ satisfies $\langle x,\theta\rangle=0$, so samples taken outside block $i^\star$ are uninformative; moreover, if the algorithm outputs $\hat x\in\mathcal{X}_j$ with $j\neq i^\star$, then $\langle \hat x,\theta\rangle=0$, which is the same value as outputting the zero vector on block $i^\star$. Therefore the $(\varepsilon,\delta)$-PAC requirement for sufficiently small absolute constant $\delta$ forces $H_2\ge c_3\cdot kd^2$ for an absolute constant $c_3>0$.

$H_1\geq \Omega(kd)$ follows directly from Theorem~\ref{thm:adaptive-lower-bound}.

% We first claim that, for the set $\mathcal{X}$, any non-adaptive algorithm has minimax sample complexity at least $\Omega\!\left(\frac{kd\log(1/\delta)}{\varepsilon^2}+\frac{kd^2}{\varepsilon^2}\right)$ for the $\varepsilon$-best arm identification problem. The term $\Omega\!\left(\frac{kd\log(1/\delta)}{\varepsilon^2}\right)$ follows directly from Theorem~\ref{thm:adaptive-lower-bound} together with the fact that $\mathrm{span}(\mathcal{X})$ has dimension $kd$. The $\Omega\!\left(\frac{kd^2}{\varepsilon^2}\right)$ term arises from the block structure of $\mathcal{X}$. Each set $\mathcal{X}_i$ is a unit $\ell_2$-ball in a $d$-dimensional coordinate subspace, and a non-adaptive algorithm does not know in advance which block contains the optimal arm. As a result, the sampling budget must be spread across all $k$ blocks, requiring on the order of $\Omega\!\left(\frac{d^2}{\varepsilon^2}\right)$ samples per block to ensure that, regardless of which $\mathcal{X}_i$ contains the best arm, the algorithm can identify an $\varepsilon$-best arm within that block upon termination. We prove this claim formally in Appendix~XYZ.

